\clearpage\section{向量运算、连续性与平移}\label{g:109:1}

\textbf{对应原课件：109，第 1--7 页。}


\paragraph{本讲的主线。}
108把输出看成集合，109进一步利用输出空间 $\R^N$ 能做加法与数乘的结构。先问：两个可行集相加后是否紧、是否闭？紧集容易，闭而无界的集会有“两个分量各跑向无穷、和却收敛”的麻烦。为控制这种现象引入渐近锥，再研究凸性如何简化无穷远结构。
\[\text{向量与集合运算}\ \longrightarrow\ \text{渐近方向与闭和}
\ \longrightarrow\ \text{凸性、凸包、超平面}\ \longrightarrow\ \text{闭凸集的衰退方向}.\]
原页67以后重新介绍渐近锥并逐项证明，不是新的同名概念。后半段加强了前面的序列工具，最终证明闭凸情形 $AS=S_\infty=\operatorname{rec}S$。


\paragraph{原页1：本讲研究什么。}本讲标题是“$\R^N$ 中的向量”，但内容远超向量计算：研究集合相加是否保持闭性、凸性，以及无界集合往哪些方向延伸。经济学中多个生产集合相加就是总生产可能集合；证明这个总集合闭，才能保证极限方案仍可行。封面给出教师、教材与授课时间，不含需要证明的命题。



\paragraph{原页2：在课程知识链中的位置。}前面的极限、连续性和对应理论，本讲都要用到。主线是：向量加法连续 $\Rightarrow$ 紧集之和紧；无穷远方向不能互相抵消 $\Rightarrow$ 闭集之和闭；凸性 $\Rightarrow$ 可用超平面描述边界。下一讲的不动点定理又需要紧性与凸性。本页的目录不是一组定理，先用它定位依赖关系。



\paragraph{原页3：向量相加：逐个坐标相加。}若 $x^1=(1,2)$、$x^2=(3,-1)$，则 $x^1+x^2=(4,1)$，$-x^1=(-1,-2)$，$x^1-x^2=(-2,3)$。上标 $1,2$ 标记两个向量，并非幂。一般写 $\sum_{k=1}^K x^k$，其第 $i$ 个坐标是 $\sum_{k=1}^K x_i^k$。课件末行的求和指标排版有混用，记住“$k$ 数向量、$i$ 数坐标”即可。零向量的每个坐标都是 $0$。



\paragraph{原页4：数乘及平行四边形直观。}$t x=(t x_1,\ldots,t x_N)$：$t>1$ 拉长，$0<t<1$ 缩短，$t<0$ 还反向，$t=0$ 变为原点。向量加法可看成先走 $x^1$ 再走 $x^2$；终点是平行四边形的对角顶点。若 $x=0$，所谓“直线 $0x$”不唯一，但代数定义仍给 $t0=0$，不受影响。



\paragraph{原页5：为什么加法和数乘连续。}\paragraph{加法的完整验证。} 若 $u_k\to u$、$v_k\to v$，则
\[\norm{(u_k+v_k)-(u+v)}\le\norm{u_k-u}+\norm{v_k-v}\longrightarrow0.\]
\paragraph{数乘的完整验证。} 若 $t_k\to t$、$x_k\to x$，把差拆成 $t_k(x_k-x)+(t_k-t)x$。因为收敛数列 $(t_k)$ 有界，存在 $M$ 使 $|t_k|\le M$，于是
\[\norm{t_kx_k-tx}\le M\norm{x_k-x}+|t_k-t|\norm x\longrightarrow0.\]
这是课件“显然连续”背后的理由。考试不能只说“极限可以交换”，应给出连续性或以上估计。



\paragraph{原页6：点与向量是两种看法。}同一个数组 $(x_1,\ldots,x_N)$，说“点”时强调位置，说“向量”时强调从原点到该点的位移。比如 $x-y$ 是从 $y$ 指向 $x$ 的位移。后面的集合平移和法向量都用这种转换；不需要另学一种新的对象。



\paragraph{原页7：平移。}固定 $a$，$T_a(x)=x+a$ 将所有点沿同一方向搬动。它连续，逆映射 $T_{-a}$ 也连续，因此保留开集、闭集、紧集、连通性。证明闭性也可直接做：若 $x_k+a\to z$ 且 $x_k\in C$，则 $x_k\to z-a$；$C$ 闭推出 $z-a\in C$，故 $z\in C+a$。平移不改变无穷远方向，见第31页。



\clearpage\section{集合和、集合差与闭包}\label{g:109:8}

\textbf{对应原课件：109，第 8--10 页。}


\paragraph{集合加法与集合差，一定把选点写出来。}
\[A+B=\{a+b:a\in A,b\in B\},\quad -A=\{-a:a\in A\},\quad A-B=A+(-B).
\]
$A-B$ 在本节是所有向量差，\textbf{不是}集合差 $A\setminus B$。例如 $A=B=\{0,1\}$，则 $A-B=\{-1,0,1\}$，而 $A\setminus B=\varnothing$。

老师图中的水平线段 $A=[0,1]\times\{0\}$ 与竖直线段 $B=\{0\}\times[0,1]$ 相加是整个正方形 $[0,1]^2$：和点是 $(a,0)+(0,b)=(a,b)$，两坐标独立选取。不是只把对应端点相加。
\paragraph{为什么“闭集之和闭”不能直接说。}
在 $\R^2$ 取闭集 $A=\{(t,1/t):t>0\}$ 与 $B=\{(-s,0):s\ge0\}$。$A$ 闭：若其中的点趋于有限 $(a,b)$，则 $a\ge0$；若 $a=0$，第二坐标会趋于无穷，与收敛矛盾，所以 $a>0$ 且 $b=1/a$。$B$ 是闭射线。
但 $z_n=(n,1/n)+(-n,0)=(0,1/n)\to(0,0)$，而任何 $A+B$ 的点第二坐标严格正，故 $(0,0)\notin A+B$。所以和不闭。两个分量的无穷远方向相反，正是后面正半独立条件排除的情形。


\paragraph{原页8：集合相加：闵可夫斯基和。}$A+B=\{a+b:a\in A,b\in B\}$，意为两边各挑一个点，再把所有可能的和收集起来。它与并集 $A\cup B$ 完全不同。例如 $[0,1]+[2,4]=[2,5]$，而并集是两段。证明 $z\in A+B$ 时，应展示分解 $z=a+b$；证明不存在这种分解即可证明不属于和集。



\paragraph{原页9：两条线段相加得到正方形。}课件中的第一条线段是 $A=\{(s,0):0\le s\le1\}$，第二条是 $B=\{(0,t):0\le t\le1\}$。任取两点相加得到 $(s,t)\in[0,1]^2$，故 $A+B\subset[0,1]^2$；反过来每个 $(s,t)$ 都可写为 $(s,0)+(0,t)$，故另一包含也成立。集合等式通常要这样证明两个方向，不能只凭图形判断。



\subsection{命题与证明：原页 10}

\textbf{命题。}$\sum_i\overline{C_i}\subseteq\overline{\sum_i C_i}$（有限和）。

\paragraph{老师的证明思路。}每个坐标点分别从原集中逼近，再对有限和取极限。

\paragraph{逐步推导。}
给有限个 $C_i\subseteq\R^N$，要证 $\sum_i\overline{C_i}\subseteq\overline{\sum_iC_i}$。若某因子空，左边空，包含直接成立。以下任取 $z\in\sum_i\overline{C_i}$。和集定义允许选 $c_i\in\overline{C_i}$ 使 $z=\sum_i c_i$。
对每个 $i$，闭包序列刻画给 $c_{i,n}\in C_i$ 且 $c_{i,n}\to c_i$。这里有限个列可以共同以 $n$ 编号，不要求它们来自同一集合。令 $z_n=\sum_i c_{i,n}$，按和定义每个 $z_n\in\sum_i C_i$。三角不等式给
\[\|z_n-z\|=\left\|\sum_i(c_{i,n}-c_i)\right\|\le\sum_i\|c_{i,n}-c_i\|\to0.\]
最后极限因为有限多个非负趋零项的和趋零。因此 $z$ 是和集内点列的极限，属于和集闭包。任意 $z$ 都如此，得到所需包含。
这不证明反向包含，也不证明原和集闭；老师两个闭包横线的位置正是结论的实质。


\clearpage\section{紧集的和：老师的两条证明}\label{g:109:11}

\textbf{对应原课件：109，第 11--14 页。}


\paragraph{两个证明各自解决什么。}
连续像法把整个和集写成加法映射的像，再调用有限紧积。序列法从和点 $z_n$ 开始，为每项取分解，再依次控制每个分量。不能从 $z_n$ 收敛倒推出分量收敛；$n+(-n)=0$ 就反驳这种推理。紧性是在\emph{分量所在集合}上使用的。


\subsection{命题与证明：原页 11、12}

\textbf{命题。}有限个紧集的和紧。

\paragraph{老师的证明思路。}保留原页的连续像或逐因子抽子列路线；闭性和有界性须分别解释。

\paragraph{逐步推导。}
设 $C_1,\ldots,C_m$ 紧。若有空因子，和集空，紧。以下均非空。
\paragraph{原页11：统一有界性}
每个紧因子有界，取 $M_i$ 使 $\|c_i\|\le M_i$ 对其全部元素成立。任意和点 $z$ 按定义可写 $z=\sum_i c_i$。三角不等式逐次应用给
\[\|z\|\le\sum_i\|c_i\|\le\sum_iM_i=:M.\]
$M$ 是事先由各因子确定的有限常数，不依赖当前 $z$，所以整个和集有界。
\paragraph{原页12：用连续像补闭性}
定义 $F:\prod_i\R^N\to\R^N$，$F(c_1,\ldots,c_m)=\sum_i c_i$。若每个坐标块趋于极限，有限和的误差满足上面相同三角估计趋零，故 $F$ 连续。有限紧积 $K=\prod_iC_i$ 紧；它在 $F$ 下的像正是集合和，因为输出恰是各因子各选一个元素相加的全部可能。
连续紧像定理给 $F(K)$ 紧，因此在有限维也闭。与前段有界性合起来，满足老师先陈述的闭有界判据；事实上连续像步骤已经直接给紧性，不必再重复推导。
绝不能把“每个 $C_i$ 闭”当作“和闭”的理由，一般闭和未必闭；原页专门用紧积连续像正是补上这项。


\subsection{命题与证明：原页 13、14}

\textbf{命题。}有限个紧集的和紧。

\paragraph{老师的证明思路。}保留原页的连续像或逐因子抽子列路线；闭性和有界性须分别解释。

\paragraph{逐步推导。}
这是老师另一条独立路线，保留全部有限次嵌套子列。任取和集内序列 $z_n$，为每项选一个分解 $z_n=\sum_{i=1}^m c_{i,n}$，其中 $c_{i,n}\in C_i$。分解可以不唯一，任取一组即可。
紧 $C_1$ 给第一坐标子列 $c_{1,n^{(1)}_k}\to c_1\in C_1$。将\emph{全部坐标及原和列}同时限制到这些下标，再用 $C_2$ 紧性从中抽子列使第二坐标收敛；第一坐标已有极限，抽子列后仍保持。继续有限 $m$ 次后得到同一最终下标 $n_k$，使所有 $c_{i,n_k}\to c_i\in C_i$。
因此
\[z_{n_k}=\sum_i c_{i,n_k}\longrightarrow\sum_i c_i=:z,\qquad z\in\sum_i C_i.\]
原页14正是这个最后有限和取极限与极限归属步骤，必须接在原页13之后，不能漏掉。因为任取和集序列都有收敛到集内的子列，度量空间序列紧判据给和集紧。
老师每一步同时保留之前各坐标极限十分重要。若独立为每个坐标抽取子列却不嵌套，可能得不到共同下标，也就无法对原和列取极限。


\clearpage\section{函数相加与对应相加}\label{g:109:15}

\textbf{对应原课件：109，第 15--17 页。}


\paragraph{函数与对应的相加差别。}
函数和 $f(x)=\sum_i f_i(x)$ 每项由 $x$ 唯一确定，逐项收敛即可相加。对应和 $\Phi(x)=\sum_i\Phi_i(x)$ 的一个 $y$ 可能有很多分解。下半连续先固定旧分解再逐项逼近；上半连续则先分解每个新输出，借紧性在所有分解里找到可取极限的一条子列。下面保留老师的这两个不同方向。


\paragraph{原页15：有限个连续函数相加仍连续。}设 $f=\sum_{k=1}^K f_k$，要证在 $x_0$ 连续。给定 $\eps>0$，每个 $f_k$ 连续，选 $\delta_k>0$ 使 $\norm{x-x_0}<\delta_k$ 时 $\norm{f_k(x)-f_k(x_0)}<\eps/K$。取 $\delta=\min_k\delta_k>0$，三角不等式给
\[\norm{f(x)-f(x_0)}\le\sum_k\norm{f_k(x)-f_k(x_0)}<\eps.\]
“有限”保证最小值仍为正；无限级数不能直接照抄，通常要增加一致收敛等条件。



\paragraph{原页16：对应相加：为什么需要控制值域。}对应 $\varphi_k(x)$ 给出一个可选集合；$\varphi(x)=\sum_k\varphi_k(x)$ 给出所有总选择。课件把各值域放在紧集 $T_k$ 中，以便从分解的各分量抽取收敛子列。采用本教材的紧值、闭图式上半连续约定时，上半连续性得以保持。若使用一般拓扑教材中不自动要求闭值的“开集定义”，须另外明确紧值等条件，不能不加条件地把闭图与上半连续等同。下半连续表示旧选择能由邻近参数下的选择逼近；求和的下半连续不需要各 $T_k$ 紧。



\subsection{命题与证明：原页 17}

\textbf{命题。}共同紧陪域约定下，有限个上半连续对应的和具有相应上半连续性。

\paragraph{老师的证明思路。}给和中的每项选择一个分解，再逐分量抽子列，不能直接假设每个分量已收敛。

\paragraph{逐步推导。}
设因子对应值域各被一个固定紧集 $T_i$ 容纳，$\Phi(x)=\sum_i\Phi_i(x)$。本节采用闭图式上半连续，紧陪域允许它与常用定义对应。
\paragraph{上半连续方向}
给 $x_n\to x_0,y_n\to y_0,y_n\in\Phi(x_n)$。逐项取分解 $y_n=\sum_i y_{i,n}$，$y_{i,n}\in\Phi_i(x_n)\subseteq T_i$。不能从和收敛直接说每个分量收敛：例如 $n+(-n)=0$。这里应使用紧 $T_i$，依次嵌套抽取有限次子列，使同一最终指标下每个 $y_{i,n_k}\to y_{i,0}$。
因子闭图式上半连续给 $y_{i,0}\in\Phi_i(x_0)$。有限和取极限使 $y_{n_k}\to\sum_i y_{i,0}$，原列收敛使此子列也趋于 $y_0$，由极限唯一 $y_0=\sum_i y_{i,0}\in\Phi(x_0)$，完成所需。
\paragraph{下半连续方向}
给 $x_n\to x_0$ 及任意 $y_0\in\Phi(x_0)$，先固定一个旧分解 $y_0=\sum_i y_{i,0}$。各因子下半连续分别给整列 $y_{i,n}\in\Phi_i(x_n)$ 趋于 $y_{i,0}$。令 $y_n=\sum_i y_{i,n}$，则可行归属与有限和极限分别给 $y_n\in\Phi(x_n)$、$y_n\to y_0$。这次不需要抽子列，因为各因子都能对给定整列输入提供逼近。
原课件上半连续段直接写出各分量的极限，省略的正是紧性与共同子列，不能把它当作代数恒等式。


\clearpage\section{线段、射线、锥与正半独立}\label{g:109:18}

\textbf{对应原课件：109，第 18--23 页。}


\paragraph{同一公式，参数范围决定几何对象。}
设 $x\ne y$，$z=x+t(y-x)$。$t\in\R$ 得直线；$t\ge0$ 得从 $x$ 出发朝 $y$ 的闭射线；$t>0$ 得去掉起点的开射线；$0\le t\le1$ 得闭线段。$x=y$ 时线段退化为一个点。

权重和1的 $t_1x+t_2y$ 是仿射组合，允许负权重；\textbf{凸组合还要求权重非负}。比如 $2x-y$ 的权重和也是1，但一般在线段外。原页18的 weighted average 只有和1条件，不能在证明中自动当成凸组合。
\paragraph{锥与顶点。}
顶点 $a$ 的锥 $C$ 满足：$a\in C$，且 $y\in C,t\ge0\Rightarrow a+t(y-a)\in C$。顶点0时写成 $y\in C,t\ge0\Rightarrow ty\in C$。锥不一定凸：两条坐标正半轴的并是锥，但 $(1,0),(0,1)$ 的中点不在里面。
\paragraph{正半独立的全体定义。}
锥 $C_1,\ldots,C_m$ 正半独立指
\[d_i\in C_i\ (i=1,\ldots,m),\quad\sum_i d_i=0
\quad\Longrightarrow\quad d_1=\cdots=d_m=0.
\]
它禁止从各锥选出的非零方向彼此抵消。两个锥时可化成 $C_1\cap(-C_2)=\{0\}$，下面逐向证明。它不同于要求所有方向向量线性独立：这里只管\emph{合法锥方向的零和}。


\paragraph{原页18：加权平均与凸组合的区别。}$t_1+t_2=1$ 时 $t_1x^1+t_2x^2$ 称仿射组合，课件称 weighted average。若再有 $t_1,t_2\ge0$，才是凸组合。比如 $2x^1-x^2$ 的系数和为 $1$，却可能跑到线段外。后面证明凸性时必须检查“系数非负”和“系数和为 $1$”两个条件。



\paragraph{原页19：直线、射线：参数的范围决定形状。}统一公式是 $z=x+t(y-x)$。$t\in\R$ 给整条直线；$t\ge0$ 给从 $x$ 出发经过 $y$ 的闭射线；$t>0$ 去掉起点，给开射线。取 $x=(1,1)$、$y=(3,2)$，则 $z=(1+2t,1+t)$。当 $t=0,1,2$ 时分别是起点、$y$、继续向前的点。“闭射线”指包含端点，并不是有界。



\paragraph{原页20：图5：方向相同，起点不同。}图中下方射线从 $0$ 出发，方向是 $y-x$；上方射线由下方整条平移 $x$ 得到，起点变成 $x$ 并经过 $y$。因此“方向”不包含起点信息。图里的斜率在二维画面可看见，但 $\R^N$ 中更可靠的描述始终是 $x+t(y-x)$。



\paragraph{原页21：线段与退化线段。}把第19页的参数限制为 $0\le t\le1$，得到 $[x^1,x^2]$。两端都包括在内。若 $x^1=x^2$，任意 $t$ 都给同一个点，线段退化成单点。以后用“包含任意两点的整段”定义凸性时，单点也应允许，不能强行要求两点不同。



\paragraph{原页22：锥：沿射线可以一直走。}顶点为 $x$ 的锥 $C$ 要满足：$x\in C$，且任意 $y\in C$、$t\ge0$ 都有 $x+t(y-x)\in C$。顶点为 $0$ 时简化为 $t y\in C$。\emph{锥不一定凸}：二维正横轴与正纵轴的并是锥，却不包含 $(1,0)$ 与 $(0,1)$ 的中点。凸锥还要容许不同方向的非负线性组合。



\subsection{命题与证明：原页 23}

\textbf{命题。}两个含0的锥正半独立当且仅当 $C_1\cap(-C_2)=\{0\}$。

\paragraph{老师的证明思路。}把 $d_1+d_2=0$ 改写为 $d_1=-d_2$，在交集中判断。

\paragraph{逐步推导。}
设两锥都含0。正半独立的定义为：若 $d_1\in C_1,d_2\in C_2$ 且 $d_1+d_2=0$，则两者均0。
若 $C_1\cap(-C_2)=\{0\}$，给这样一组相加为0的方向，有 $d_1=-d_2$，所以 $d_1$ 同时在 $C_1$ 和 $-C_2$ 中，必为0，继而 $d_2=-d_1=0$，正半独立成立。
反向假设正半独立，任取 $d\in C_1\cap(-C_2)$，则 $d\in C_1$ 且 $-d\in C_2$，相加为0；定义给 $d=0$，所以交集包含于 $\{0\}$。两锥都含0又使0属于交，因此交恰为 $\{0\}$。
负锥 $-C_2=\{-c:c\in C_2\}$ 不是补集；交集中的点表示一个方向既由第一锥向前走，也由第二锥反向走，正是会互相抵消的情况。


\clearpage\section{三种无穷远描述：先把定义分开}\label{g:109:24}

\textbf{对应原课件：109，第 24--28 页。}


\paragraph{先列三种定义，后文再证明它们何时相同。}
设 $C\ne\varnothing$。截尾集 $C^M=\{x\in C:\|x\|\ge M\}$ 只保留离原点很远的点；$\Gamma(C^M)$ 是包含它的最小闭锥，不要求凸。定义
\[AC=\bigcap_{M\ge0}\Gamma(C^M).
\]
因为阈值越大截尾越小，所以在每个阈值下仍保留下来的才是无穷远方向。空截尾按含0锥约定生成 $\{0\}$。

序列版本定义
\[C_\infty=\{d:\exists x_k\in C,\ \exists t_k\to+\infty,\ x_k/t_k\to d\}.
\]
标度 $t_k$ 可以依赖所求方向，且不要求恰等于 $\|x_k\|$。$d=0$ 也合法：固定 $a\in C$，令 $x_k=a,t_k=k$。$d\ne0$ 才迫使 $\|x_k\|\to\infty$。

衰退锥定义
\[\operatorname{rec}C=0^+C=\{d:\forall a\in C\ \forall t\ge0,\ a+td\in C\}.
\]
它要求从\emph{每个基点}出发都能沿此方向一直留在集合里，比仅仅能挑某条趋远点列强。一般只有 $\operatorname{rec}C\subseteq C_\infty$；非空\textbf{闭凸}时才等号。第27页与第75页末的 $AC=0^+C$ 必须连同这个条件读。
\paragraph{三个具体例子。}
$C=[0,1]$ 有 $C_\infty=\{0\}$。$C=[1,\infty)$ 有 $C_\infty=[0,\infty)$：非负 $d$ 可取 $x_k=1+kd$，负极限不可能由非负的比值得到。$C=\{(t,t^2):t\ge0\}$ 的渐近锥是竖直向上的射线：取 $t\to\infty$ 并以 $t^2$ 为标度得 $(0,1)$；但这条抛物线上不能沿竖直方向走直线，所以衰退锥仅为 $\{0\}$。这说明“趋远方向”不等于“集合里已有射线”。


\paragraph{原页24：渐近锥：只保留越来越远的点。}先截取 $S^M=\{x\in S:\norm x\ge M\}$，再作包含它的最小闭锥 $\Gamma(S^M)$，最后取所有 $M\ge0$ 的交。这样的方向必须在无论多远的截尾中都能被逼近。锥在这里总含原点，特别约定 $\Gamma(\varnothing)=\{0\}$。每个 $\Gamma(S^M)$ 闭且对非负数乘封闭，交仍闭且是锥，所以 $AS$ 是闭锥。它不一定凸，除非再对 $S$ 加凸性条件。



\paragraph{原页26：序列定义：把远点缩小后取极限。}对非空 $C$，$d\in C_\infty$ 等价于存在 $x_k\in C$、$t_k\to+\infty$ 使 $x_k/t_k\to d$。这里可以删掉有限个初项，使 $t_k>0$。例：$C=\{(s,1):s\ge0\}$ 的渐近锥为 $\{(a,0):a\ge0\}$，因为缩放后第二坐标 $1/t_k\to0$。零方向总存在：固定一点 $x\in C$，用 $x/k\to0$。非零方向才表示真正逃向无穷远。



\paragraph{原页27：衰退锥：从每个点都能沿它走。}定义 $0^+C=\{d:x+\lambda d\in C\text{ 对所有 }x\in C,\lambda\ge0\}$。这比“存在远点列沿该方向”强得多。一般只有 $0^+C\subset C_\infty$；$C$ 非空闭凸时才可保证相等，完整证明见第110--111页。\textbf{条件提示：} 本页末行 $AC=0^+C$ 若只假设凸而未假设闭，并不普遍成立；第112页正好给反例。不要把不同作者的术语相同误当作无条件等价。



\paragraph{原页28：渐近锥性质总览与两个边界条件。}正确读法：$C_\infty$ 闭且是锥；$(\overline C)_\infty=C_\infty$；若 $C$ 本身是锥，则 $C_\infty=\overline C$，横线不能漏。证明在第77--80页。
归一化方向集 $C_N$ 是无界点列 $x_k$ 的单位方向 $x_k/\norm{x_k}$ 的极限集。为兼容有界 $C$ 导致 $C_N=\varnothing$ 的情形，等式应写
\[C_\infty=\{0\}\cup\{\lambda u:\lambda>0,\ u\in C_N\}.\]
“$C$ 有界当且仅当 $C_\infty=\{0\}$”按本页上下文要求 $C$ 非空，空集另行处理。



\subsection{命题与证明：原页 25}

\textbf{命题。}截尾集合及其生成闭锥随截尾阈值增大而递减；生成闭锥须保留闭包。

\paragraph{老师的证明思路。}先比较原截尾集包含，再用生成闭锥的单调性取交。

\paragraph{逐步推导。}
定义截尾集 $S^M=\{x\in S:\|x\|\ge M\}$。若 $M_1\le M_2$，任何满足更强条件 $\|x\|\ge M_2$ 的点也满足较弱条件，所以 $S^{M_2}\subseteq S^{M_1}$。
令 $\Gamma(E)$ 为包含 $E$ 且含0的最小闭锥。若 $E\subseteq F$，每个含 $F$ 的闭锥也含 $E$；尤其 $\Gamma(F)$ 是 $E$ 的一个闭锥超集，由最小性 $\Gamma(E)\subseteq\Gamma(F)$。应用于截尾集，得 $\Gamma(S^{M_2})\subseteq\Gamma(S^{M_1})$。所以这族闭锥随阈值递减，其集合列极限是交 $A S=\bigcap_M\Gamma(S^M)$。
老师再写生成锥的公式时必须保留闭包：
\[\Gamma(E)=\overline{\{\alpha x:\alpha\ge0,x\in E\}}\]
对非空 $E$ 成立。理由是未闭包的非负倍乘集已为锥且含 $E$；其闭包仍为锥，因为固定倍乘连续；任意闭锥超集又必须包含这些倍乘及其闭包，所以它确为最小。若 $E$ 空，则按“含0”的约定最小锥为 $\{0\}$，须单独处理。不能把生成闭锥误当凸包，也不能漏闭包。


\clearpage\section{有界判据、平移不变与乘积包含}\label{g:109:29}

\textbf{对应原课件：109，第 29--33 页。}


\paragraph{有界性为何要看单位向量。}
想从无界构造\emph{非零}渐近方向，不能随便取标度 $t_n$：太大可能把所有点压成0。老师取 $t_n=\|x_n\|$，让比值范数恒1；单位球面紧给子列，而极限范数仍1，保证方向不是0。这是本讲反复出现的标准构造。
\paragraph{先记住只成立的包含。}
\[(C_1\times\cdots\times C_m)_\infty
\subseteq(C_1)_\infty\times\cdots\times(C_m)_\infty.
\]
左侧所有坐标来自同一个趋无穷标度；右侧各因子可以各有一条不同标度。证明反向时若直接把不同标度拼成一个，就偷换了量词。凸性带来的渐近正则性将在后面解决某些统一尺度问题。


\subsection{命题与证明：原页 29}

\textbf{命题。}非空集合有界当且仅当其序列式渐近锥只有零向量。

\paragraph{老师的证明思路。}有界时除以趋无穷标度；无界时选范数趋无穷点列并单位化抽子列。

\paragraph{逐步推导。}
本条 $C\ne\varnothing$。若 $C$ 有界，取 $M$ 使任意 $x\in C$ 有 $\|x\|\le M$。任取 $d\in C_\infty$，定义给 $x_n\in C,t_n\to+\infty$ 且 $x_n/t_n\to d$。充分大 $t_n>0$，有 $\|x_n/t_n\|\le M/t_n\to0$，所以 $d=0$。反向0确实在锥中：固定 $x\in C$，取 $x_n=x,t_n=n$，则 $x/n\to0$。故锥恰为 $\{0\}$。
反向假设锥只有0，若 $C$ 无界，可对每个 $n$ 选 $x_n\in C$ 使 $\|x_n\|>n$。令 $u_n=x_n/\|x_n\|$，范数恒1。有限维单位球面紧，存在子列 $u_{n_k}\to u$，范数连续给 $\|u\|=1$，所以 $u\ne0$。
取对应 $t_k=\|x_{n_k}\|\to\infty$，则 $x_{n_k}/t_k=u_{n_k}\to u$，定义使 $u\in C_\infty$，与只有0矛盾。因此 $C$ 有界。不能说任意无界列范数都趋无穷，第一步的逐层选点专门保证这一点。


\subsection{命题与证明：原页 30}

\textbf{命题。}渐近锥随原集合包含而单调；给非空集合加另一非空集合使渐近锥不缩小。

\paragraph{老师的证明思路。}先用截尾包含证明单调，再固定一个平移点应用平移不变性。

\paragraph{逐步推导。}
\paragraph{单调性}若 $S_1\subseteq S_2$，则每个阈值 $M$ 下 $S_1^M\subseteq S_2^M$。生成闭锥的单调性给 $\Gamma(S_1^M)\subseteq\Gamma(S_2^M)$。任取 $d\in AS_1$，它属于左侧每一项，因此属于右侧每一项，即属于 $AS_2$。故 $AS_1\subseteq AS_2$。
\paragraph{非空和集推论}设 $T\ne\varnothing$，固定一个 $a\in T$，则 $S+\{a\}\subseteq S+T$。先用下一条已证明的平移不变性 $A(S+\{a\})=AS$，再用刚证单调性，得到
\[AS=A(S+\{a\})\subseteq A(S+T).\]
若不假设 $T$ 非空，就无法选 $a$，而和集可能空，因此非空性确实参与论证。引用平移不变性时需把它视为本组随后独立证明的引理，不能拿本推论反过来证明它。


\subsection{命题与证明：原页 31}

\textbf{命题。}平移不改变渐近锥：$(C+a)_\infty=C_\infty$。

\paragraph{老师的证明思路。}把见证序列加减固定向量；除以趋无穷标度后平移项趋零。

\paragraph{逐步推导。}
任取 $d\in(C+a)_\infty$。按定义有 $z_n\in C+a,t_n\to\infty$ 且 $z_n/t_n\to d$。逐项写 $z_n=x_n+a$，其中 $x_n\in C$。则
\[\left\|\frac{x_n}{t_n}-d\right\|\le\left\|\frac{z_n}{t_n}-d\right\|+\frac{\|a\|}{t_n}\to0.\]
两项趋零的理由分别是见证序列的已知收敛、固定平移除以趋无穷标度。因此 $d\in C_\infty$，得到 $(C+a)_\infty\subseteq C_\infty$。
反向对任意 $d\in C_\infty$，取原见证 $x_n/t_n\to d$，令 $z_n=x_n+a\in C+a$。同一误差估计给 $z_n/t_n\to d$，故 $d\in(C+a)_\infty$。或者按老师写法，将第一包含中的集合换为 $C+a$、平移换为 $-a$，也得到反向。两包含合成等号。
这说明固定原点位置的改变被无穷尺度消掉，不是说原集合本身平移后相等。


\subsection{命题与证明：原页 32}

\textbf{命题。}平移不改变渐近锥：$(C+a)_\infty=C_\infty$。

\paragraph{老师的证明思路。}把见证序列加减固定向量；除以趋无穷标度后平移项趋零。

\paragraph{逐步推导。}
老师的等价链必须展开为见证序列的双向转换：
\begin{align*}
d\in C_\infty
&\iff\exists x_n\in C,\ t_n\to\infty:\ x_n/t_n\to d\\
&\iff\exists z_n=x_n+a\in C+a,\ t_n\to\infty:\ z_n/t_n\to d\\
&\iff d\in(C+a)_\infty.
\end{align*}
中间等价不是仅靠重命名，因为表达式多出了 $a/t_n$。正向若 $x_n/t_n\to d$，则 $z_n/t_n=x_n/t_n+a/t_n\to d+0=d$；反向若 $z_n/t_n\to d$，则 $x_n=z_n-a\in C$ 且 $x_n/t_n=z_n/t_n-a/t_n\to d$。$a$ 固定保证 $a/t_n\to0$，标度保持不变，所以两方向见证都合法。
最后截尾定义与序列定义已经证明等价时，可将 $C_\infty$ 换写为 $AC$，得到原页的记号。不能只对其中一个包含验证见证转换便称为等价链。


\subsection{命题与证明：原页 33}

\textbf{命题。}有限积的渐近锥包含于各因子渐近锥的积。

\paragraph{老师的证明思路。}对共同标度下的向量极限逐坐标取极限；原页截尾等式的错误另外说明。

\paragraph{逐步推导。}
设 $C=\prod_{i=1}^mC_i$。若积为空，包含按相应空锥约定单独成立；以下取 $d\in C_\infty$，由定义选 $x_n=(x_{1n},\ldots,x_{mn})\in C$ 与同一 $t_n\to\infty$，满足 $x_n/t_n\to d=(d_1,\ldots,d_m)$。
对每个固定 $i$，积归属给 $x_{in}\in C_i$，坐标投影连续给 $x_{in}/t_n\to d_i$。这正是 $d_i\in(C_i)_\infty$ 的完整见证，故全部坐标合起来得到
\[d\in\prod_i(C_i)_\infty.\]
由任意 $d$ 得 $C_\infty\subseteq\prod_i(C_i)_\infty$。0方向也被此推理覆盖，无须跳过。
反向不能从各自见证直接拼接，因为不同因子的 $t_n$ 可能不相同；本条只证明包含。原页89另用截尾的证明存在错误，见其单列辨错步骤。


\clearpage\section{渐近锥怎样控制交集的有界性}\label{g:109:34}

\textbf{对应原课件：109，第 34--36 页。}


\paragraph{交有界与和闭是两种不同问题。}
若各集合的渐近锥交只有0，则共同集合不能有任意远的点，因而交有界。若各集合的渐近锥正半独立，则加法中不能发生无穷远抵消，因而在\emph{各集合还闭}时，和闭。前一个条件是共同方向，后一个条件是方向之和为0，不能互相替代。


\paragraph{原页34：两个核心应用：交有界，和闭。}交集有界的充分条件是所有渐近锥的交仅为 $\{0\}$；闭集之和闭的充分条件是它们的渐近锥正半独立。两者都在阻止某种“逃到无穷”的行为：前者阻止同时在所有集合中逃走，后者阻止各分量逃走但总和相互抵消。它们是\emph{充分条件}，不要写成必要条件。例如整个 $\R$ 加自己仍闭，两个渐近锥却能抵消。



\subsection{命题与证明：原页 35、36}

\textbf{命题。}若各集合渐近锥的交至多为零方向，则原集合的交有界。

\paragraph{老师的证明思路。}交集包含于每一因子，先用锥单调性，再用有界性判据。

\paragraph{逐步推导。}
设 $C=\bigcap_\gamma C_\gamma$，且 $\bigcap_\gamma(C_\gamma)_\infty=\{0\}$。若 $C$ 为空，空集本来有界；以下非空。
固定任意 $\gamma$，有 $C\subseteq C_\gamma$。渐近锥单调性给 $C_\infty\subseteq(C_\gamma)_\infty$。对所有 $\gamma$ 都成立，故
\[C_\infty\subseteq\bigcap_\gamma(C_\gamma)_\infty=\{0\}.\]
因为 $C$ 非空，固定一点并除以 $n$ 可证 $0\in C_\infty$。因此不是仅有包含，而是 $C_\infty=\{0\}$。由前条“非空集合有界当且仅渐近锥只有0”，得到 $C$ 有界。
如果按老师另一版写成截尾，零以外没有渐近方向便意味着不可能存在任意远的交集点；否则将其归一化，单位球面紧性会给一个非零方向，落在上述所有锥中，矛盾。不能直接从交锥只有0无理由写“某阈值后截尾为空”，中间正是这个有界性判据。


\clearpage\section{线性像定理：归一化反证的完整链条}\label{g:109:37}

\textbf{对应原课件：109，第 37--41 页。}


\paragraph{完整命题及记号。}
设非空闭集 $C\subseteq\R^N$，线性映射 $L:\R^N\to\R^K$，且
\[\ker L\cap C_\infty=\{0\},\qquad
\ker L=L^{-1}(\{0\})=\{u:Lu=0\}.
\]
则 $L(C)$ 闭，并且 $L(C_\infty)=(L(C))_\infty$。
核是被映射压成0的方向；条件说集合不能沿某个被压扁的非零方向逃向无穷。$L$ 连续本身只保证紧像紧，\textbf{不保证任意闭像闭}。
\paragraph{老师四页证明的任务表。}
第38页：像有界但原像无界会产生核中的非零渐近方向。第39页：由原像有界抽子列，用闭性和连续性证像闭。第40页：给方向的见证取线性像，得容易的包含。第41页：给像的见证找原像，并\emph{重新证明缩放原像有界}，得难的包含。特别注意 $x_n/t_n$ 未必属于 $C$，不能把第38页直接套到它上面。

\paragraph{投影为什么可能把闭集变成不闭的集？}
仍取闭双曲线支 $C=\{(t,1/t):t>0\}$，令 $L(u,v)=u$。像为 $(0,\infty)$，不闭。问题出在输入 $(1/n,n)$ 沿竖直方向逃远，而像 $1/n\to0$；将它除以 $n$ 得 $(1/n^2,1)\to(0,1)$。因此 $(0,1)\in C_\infty$，又 $L(0,1)=0$，恰好违反核条件。这个例子把抽象的“不可见的趋远方向”变成了一个具体向量。


\paragraph{原页37：线性映射何时把闭集送到闭集。}线性映射 $L$ 满足 $L(\alpha x+\beta y)=\alpha L(x)+\beta L(y)$；有限维时它连续。核 $\ker L=L^{-1}(\{0\})$ 是所有被压到零的方向。条件
\[\ker L\cap C_\infty=\{0\}\]
表示 $C$ 的非零逃逸方向不能被 $L$ 完全压扁。若 $C$ 闭，则它保证 $L(C)$ 闭，而且 $L(C_\infty)=(L(C))_\infty$。无条件成立的方向只有 $L(C_\infty)\subset(L(C))_\infty$。下面四页分别处理“防逃逸、闭性、正向包含、反向包含”。



\subsection{命题与证明：原页 38、39、40、41}

\textbf{命题。}闭 $C$ 满足 $\ker L\cap C_\infty=\{0\}$ 时，线性像 $L(C)$ 闭且 $L(C_\infty)=(L(C))_\infty$。

\paragraph{老师的证明思路。}归一化反证原像有界；抽子列证像闭；两向分别构造渐近方向。

\paragraph{逐步推导。}
设闭 $C\subseteq\R^N$，线性 $L:\R^N\to\R^K$ 满足 $\ker L\cap C_\infty=\{0\}$。有限维线性映射连续；例如矩阵系数的有限和给 $\|Lz\|\le M\|z\|$。
\paragraph{原页38/93：像有界迫使原像有界}
若有 $x_n\in C$ 的线性像有界，而原像列无界，取子列（重记）使 $r_n=\|x_n\|\to\infty$。单位向量 $u_n=x_n/r_n$ 再抽子列趋于单位 $u$。由渐近锥定义 $u\in C_\infty$。又
\[Lu=\lim_n L(x_n/r_n)=\lim_n(Lx_n)/r_n=0,\]
因为分子有界、分母趋无穷。故 $u\in\ker L\cap C_\infty$ 却 $\|u\|=1$，矛盾。得到所需有界性。
\paragraph{原页39/94：像闭}
取任意 $y_n\in L(C)$ 且 $y_n\to y$。每项选 $x_n\in C$ 使 $Lx_n=y_n$。收敛像列有界，上段使 $x_n$ 有界，故抽 $x_{n_k}\to x$。$C$ 闭给 $x\in C$，$L$ 连续给 $Lx=\lim_kLx_{n_k}=\lim_ky_{n_k}=y$。于是 $y\in L(C)$，像闭。原像不要求唯一，任意分解都受同一核条件约束。
\paragraph{原页40/95：容易方向 $L(C_\infty)\subseteq(L(C))_\infty$}
给 $d\in C_\infty$，取 $x_n\in C,t_n\to\infty,x_n/t_n\to d$。令 $y_n=Lx_n\in L(C)$，则 $y_n/t_n=L(x_n/t_n)\to Ld$，所以 $Ld$ 是像集的渐近方向。此方向不需要闭性及核条件。
\paragraph{原页41/96：反向，补上不能直接套用的缩放有界性}
任取 $y\in(L(C))_\infty$，选 $y_n\in L(C),t_n\to\infty,y_n/t_n\to y$，再选 $x_n\in C$ 使 $Lx_n=y_n$。需证 $x_n/t_n$ 有界；因为这些缩放点\emph{未必在 $C$ 中}，不能直接应用第一段。
反设缩放列无界，取子列使 $q_n=\|x_n\|/t_n\to\infty$，并进一步使 $u_n=x_n/\|x_n\|\to u$，$\|u\|=1$。这时 $\|x_n\|=q_nt_n\to\infty$，所以 $u\in C_\infty$，且
\[Lu=\lim_n\frac{Lx_n}{\|x_n\|}=\lim_n\frac{y_n/t_n}{q_n}=0.\]
又与核条件矛盾。故 $x_n/t_n$ 有界，可抽子列趋于 $d$。同一原点及标度满足定义，$d\in C_\infty$；连续性给 $Ld=\lim_nL(x_n/t_n)=y$。所以 $y\in L(C_\infty)$。两方向合成等号。原页略过的缩放有界证明是本定理最关键的补步。


\clearpage\section{闭集相加：正半独立排除抵消}\label{g:109:42}

\textbf{对应原课件：109，第 42--45 页。}


\paragraph{把加法看作线性映射，为什么有用？}
把所有分量暂时并排组成一个大向量 $u=(x_1,\ldots,x_m)$，其合法集合是闭积 $C=\prod_iC_i$。定义 $L(u)=\sum_i x_i$，目标和集就是 $L(C)$。正半独立恰好检查了线性像定理要求的核条件，因此可以把刚证明的定理完整用上。第45页是相同推理的量词展开版本；下面合并保留全部核条件验证。


\paragraph{原页42：闭集和与正半独立。}正半独立指若 $d_i\in(C_i)_\infty$ 且 $\sum_i d_i=0$，则每个 $d_i=0$。设积集 $C=\prod_iC_i$，线性映射 $L(x_1,\ldots,x_m)=\sum_i x_i$。逐坐标极限给 $C_\infty\subseteq\prod_i(C_i)_\infty$，所以核条件成立。前面的线性像定理给 $L(C)=\sum_iC_i$ 闭，且 $(\sum_i C_i)_\infty=L(C_\infty)\subseteq\sum_i(C_i)_\infty$。这说明渐近方向不能互相抵消是保证闭和的重要条件；闭集的和一般不自动闭。



\subsection{命题与证明：原页 43、44、45}

\textbf{命题。}闭集族渐近锥正半独立时，有限集合和闭，且和的渐近锥包含于各锥之和。

\paragraph{老师的证明思路。}将和表示为积集在线性求和映射下的像，检查其核与渐近锥交。

\paragraph{逐步推导。}
假设各闭 $C_i\subseteq\R^N$ 的渐近锥正半独立：若 $d_i\in(C_i)_\infty$ 且 $\sum_i d_i=0$，则每个 $d_i=0$。若某因子空，和为空，闭性直接成立；以下非空。
令 $C=\prod_{i=1}^mC_i$，并定义线性映射 $L(x_1,\ldots,x_m)=\sum_i x_i$。有限闭积闭，且 $L(C)=\sum_iC_i$。积锥包含已给
\[C_\infty\subseteq\prod_i(C_i)_\infty.\]
若 $d=(d_1,\ldots,d_m)\in\ker L\cap C_\infty$，则每个 $d_i\in(C_i)_\infty$，同时 $\sum_id_i=Ld=0$。正半独立性强迫全部 $d_i=0$，故 $d=0$。反向0确实在该交，所以核条件正好成立。应用线性像定理，$L(C)=\sum_iC_i$ 闭。
再证明老师随后的锥包含：线性像定理给 $(L(C))_\infty=L(C_\infty)$；积锥包含经 $L$ 的单调取像给
\[L(C_\infty)\subseteq L\left(\prod_i(C_i)_\infty\right)=\sum_i(C_i)_\infty.\]
所以 $(\sum_iC_i)_\infty\subseteq\sum_i(C_i)_\infty$。最后一个等号只是求和映射的像定义，不要求在一般情形积锥等于锥的积。核条件须先验证再用定理，不能仅说 $L$ 连续便断言闭像。


\clearpage\section{凸集的定义、运算与连通性}\label{g:109:46}

\textbf{对应原课件：109，第 46--54 页。}


\paragraph{凸集的定义。}
$C$ 凸指对任意 $x,y\in C$ 和任意 $t\in[0,1]$，$(1-t)x+ty\in C$。它要求\emph{整个线段}留在集合中。圆盘凸，圆周不凸；空集与单点按定义也凸。凸性不包含闭性或有界性：开球凸但不闭，整个 $\R^N$ 凸但无界。
\paragraph{原页47：多面锥。}
由有限条从0出发的闭射线相加得到
\[C=\left\{\sum_{i=1}^m\lambda_i v_i:\lambda_i\ge0\right\}.
\]
这叫凸多面锥。对两组非负系数做非负凸组合，系数仍非负，所以它凸；把所有系数同乘非负数，仍合法，所以它是锥。与凸包区别是这里\emph{不要求系数和为1}。
\paragraph{原页48的五条结论。}
任意凸集的交凸；有限凸集的和与积凸；凸集的闭包凸；凸集连通；闭凸且含0的集合包含自己的渐近锥。后两条不是“看图明显”：连通性要借线段路径和区间连通，包含渐近锥要先用凸性缩回再用闭性取极限。下面逐条证明。

\begin{center}
\begin{tikzpicture}[scale=.95]
\begin{scope}
\draw[->](-.3,0)--(3,0)node[right]{$x_1$};
\draw[->](0,-.3)--(0,2.8)node[above]{$x_2$};
\draw[very thick,blue,->](0,0)--(2.8,0);
\draw[very thick,blue,->](0,0)--(0,2.5);
\draw[dashed](0,2)--(2,0);
\fill[red](1,1)circle(2.5pt);
\node at(1.1,-.8){两射线的并：锥，但不凸};
\end{scope}
\begin{scope}[xshift=6.5cm]
\fill[blue!12](0,0)rectangle(2.7,2.5);
\draw[->](-.3,0)--(3,0)node[right]{$x_1$};
\draw[->](0,-.3)--(0,2.8)node[above]{$x_2$};
\draw[very thick,blue,->](0,0)--(2.8,0);
\draw[very thick,blue,->](0,0)--(0,2.5);
\draw[dashed](0,2)--(2,0);\fill[red](1,1)circle(2.5pt);
\node at(1.1,-.8){两射线的和：整个非负象限};
\end{scope}
\end{tikzpicture}
\end{center}
红点是 $(2,0)$ 与 $(0,2)$ 的中点。左图不包含它，所以不凸；右图包含所有这样的非负坐标点，所以凸。右侧阴影只画出非负象限的一部分，集合向右、向上无限延伸。


\paragraph{原页46：凸集与多面锥。}凸指任取 $x,y\in C$ 及 $0\le t\le1$，$tx+(1-t)y\in C$。这是整条线段都在集合中；不是只检验中点。多面锥由有限向量的非负线性组合构成，权重不要求和为1；凸组合则要求和为1。二者必须区分。



\paragraph{原页48：凸性的基本运算及连通性。}交：两端点属于每个凸因子，线段逐个属于它们，故属于交。和：写 $x=\sum x_i,y=\sum y_i$，则 $tx+(1-t)y=\sum(tx_i+(1-t)y_i)$ 属于和；积逐坐标同理。闭包：选 $x_n,y_n\in C$ 趋于闭包两点，凸组合在 $C$，极限在闭包。凸集连通：若被相对开集分离，连接两侧点的连续路径 $t\mapsto tx+(1-t)y$ 把分离拉回 $[0,1]$，与区间连通矛盾。最后若 $C$ 闭凸且含0，$x_n/t_n\to d$ 且 $t_n\to\infty$，充分大时 $x_n/t_n=(1-1/t_n)0+x_n/t_n\in C$，闭性给 $d\in C$。



\subsection{命题与证明：原页 49}

\textbf{命题。}任意凸集族的交凸。

\paragraph{老师的证明思路。}任取两端点，凸组合在每个因子中，因而在交中。

\paragraph{逐步推导。}
设 $C=\bigcap_\gamma C_\gamma$，每个 $C_\gamma$ 凸。若交为空，凸性条件“任取交中两点”无对象，按定义空集凸。否则任取 $x,y\in C$ 及 $t\in[0,1]$。交定义使 $x,y$ 同时在每个 $C_\gamma$；对固定任意 $\gamma$，其凸性给 $tx+(1-t)y\in C_\gamma$。对所有 $\gamma$ 成立，所以该组合在交 $C$ 中。
已经验证每一对端点和每个权重，故 $C$ 凸。不必假设交闭，也不必假设指标族有限；这是纯粹逐因子全称条件的证明。


\subsection{命题与证明：原页 50}

\textbf{命题。}有限个凸集的和与积均凸。

\paragraph{老师的证明思路。}分解两端点后，以同一个权重逐因子作凸组合。

\paragraph{逐步推导。}
\paragraph{和凸}设 $C_i$ 凸，任取 $x,y\in\sum_iC_i$，分别选分解 $x=\sum_ix_i,y=\sum_iy_i$，其中 $x_i,y_i\in C_i$。给任意 $t\in[0,1]$，各因子凸使 $z_i=tx_i+(1-t)y_i\in C_i$。有限加法分配给
\[tx+(1-t)y=t\sum_ix_i+(1-t)\sum_iy_i=\sum_i[tx_i+(1-t)y_i]=\sum_iz_i,\]
故它是合法和点。对所有端点与权重成立，和凸。
\paragraph{积凸}任取 $x=(x_i),y=(y_i)\in\prod_iC_i$。向量逐坐标运算给 $tx+(1-t)y=(tx_i+(1-t)y_i)_i$。每个坐标在对应凸因子中，所以整个向量在积中，积凸。
同一个权重 $t$ 必须用于全部坐标或分量，因为我们要证明的是\emph{原两个完整向量之间的线段}，不是分别选不相关权重得到的另一点。若某因子空，和或积为空，凸性按定义成立。


\subsection{命题与证明：原页 51}

\textbf{命题。}凸集的闭包凸。

\paragraph{老师的证明思路。}用原集内序列逼近两端点，对对应凸组合取极限。

\paragraph{逐步推导。}
设 $C$ 凸。任取 $x,y\in\overline C$ 和 $t\in[0,1]$。由闭包序列刻画，分别取 $x_n,y_n\in C$ 使 $x_n\to x,y_n\to y$。因 $C$ 凸，$z_n=tx_n+(1-t)y_n\in C$。误差估计
\[\|z_n-[tx+(1-t)y]\|\le t\|x_n-x\|+(1-t)\|y_n-y\|\to0\]
表明原集中的该点列趋于所需组合。闭包包含原集内序列的全部极限，所以 $tx+(1-t)y\in\overline C$，完成定义。
原课件开头若写“$x,y$ 在凸集 $C$ 中”，应明确这里选择的是\emph{闭包}中的端点；否则只证明原集凸，尚未涉及闭包。权重固定再取极限，不能在过程中让权重未经控制地随 $n$ 改变。


\subsection{命题与证明：原页 52}

\textbf{命题。}凸集连通。

\paragraph{老师的证明思路。}反设可分离，把连接两侧点的线段参数拉回区间，违反区间连通。

\paragraph{逐步推导。}
设凸集 $C$ 不连通，反设有非空不交相对开集 $U,V$ 并为 $C$。选 $x\in U,y\in V$。令连续路径 $h:[0,1]\to\R^N$ 为 $h(t)=(1-t)x+ty$；凸性保证每个路径点都在 $C$，所以可看作到 $C$ 的连续函数。
令 $I_U=h^{-1}(U),I_V=h^{-1}(V)$。相对开集的连续逆像在 $[0,1]$ 中相对开。因 $h(0)=x\in U,h(1)=y\in V$，两逆像均非空；因 $U,V$ 不交，逆像不交；因每个路径点在 $C=U\cup V$，它们又覆盖整个 $[0,1]$。
于是区间被两个非空不交相对开集分离，与实区间连通矛盾。故 $C$ 连通。证明用到凸性是保证\emph{整条}连接路径在 $C$，仅知道端点在 $C$ 不能保证这一点。


\subsection{命题与证明：原页 53}

\textbf{命题。}含0的闭凸集包含其渐近锥。

\paragraph{老师的证明思路。}将渐近方向见证点缩小到0与该点的线段上，最后用闭性。

\paragraph{逐步推导。}
老师这一版使用截尾生成闭锥，须补足闭包逼近。设闭凸 $C$ 含0，$d\in AC$。若 $d=0$，已在 $C$。否则选足够大的阈值整数 $M$。由 $d\in\Gamma(C^M)$，可取 $z_M\in C$、$\|z_M\|\ge M$ 和 $\alpha_M\ge0$ 使
\[\|\alpha_Mz_M-d\|<1/M.\]
当 $M>\|d\|+1$ 时，$\alpha_M\|z_M\|\le\|d\|+1/M$，所以
\[0\le\alpha_M\le(\|d\|+1/M)/M<1.\]
因此 $\alpha_Mz_M=(1-\alpha_M)0+\alpha_Mz_M$ 是两个 $C$ 中端点的凸组合，在 $C$ 中。又选点误差使这些组合趋于 $d$，闭性给 $d\in C$。
原页直接说存在一个缩放点恰等于 $d$，忽略了生成\emph{闭}锥可能只含缩放极限；加上逼近与最后闭性才能严谨。保留的核心思路仍是“沿0与远点之间线段缩回”，并没有换成别的分离定理。


\subsection{命题与证明：原页 54}

\textbf{命题。}含0的闭凸集包含其渐近锥。

\paragraph{老师的证明思路。}将渐近方向见证点缩小到0与该点的线段上，最后用闭性。

\paragraph{逐步推导。}
本版改用序列式定义：给 $d\in C_\infty$，取 $z_n\in C,t_n\to+\infty$ 且 $z_n/t_n\to d$。因标度趋无穷，存在 $N$ 使 $n\ge N$ 时 $t_n\ge1$，所以 $\lambda_n=1/t_n\in[0,1]$。含0与凸性给
\[\frac{z_n}{t_n}=(1-\lambda_n)0+\lambda_nz_n\in C\qquad(n\ge N).\]
这条尾部点列全在 $C$ 中，且趋于 $d$，闭性使 $d\in C$。任取方向都如此，得到 $C_\infty\subseteq C$。
老师写“对每个 $n$ 组合都在 $C$”时，需先删去可能 $t_n<1$ 的有限项；极限不受删去有限项影响，但凸组合的合法范围不能无条件省略。含0、凸、闭分别用在端点、线段、最后极限三处。


\clearpage\section{凸包、闭凸包与凸包和公式}\label{g:109:55}

\textbf{对应原课件：109，第 55--59 页。}


\paragraph{凸包两个定义为什么一样：补全基础。}
设 $H$ 是 $S$ 的所有有限凸组合组成的集合。单个点取权重1，所以 $S\subseteq H$。若 $u=\sum_i\alpha_i s_i,v=\sum_j\beta_j t_j$ 在 $H$，则 $\theta u+(1-\theta)v$ 的新权重为 $\theta\alpha_i$ 与 $(1-\theta)\beta_j$，均非负且总和为1，所以 $H$ 凸。
另一方面，任意含 $S$ 的凸集包含 $S$ 的全部有限凸组合：两项来自定义，多项可以把前 $r-1$ 项先按总权重归一化，再与最后一项做两点凸组合，归纳成立。因此 $H$ 包含于每个含 $S$ 的凸集，又自身是这样的凸集，故
\[\conv S=\bigcap\{K:S\subseteq K,\ K\text{ 凸}\}=H.
\]
老师用上点或 $\mathrm{CH}$ 表示凸包，本文用 $\conv$；闭凸包是 $\overline{\conv S}$。它闭而且凸，并且任何含 $S$ 的闭凸集先含 $\conv S$，再因闭性含其闭包，所以它是最小闭凸超集。

\paragraph{第57--59页为何要乘积权重。}
如果 $u=\alpha a_1+(1-\alpha)a_2$、$v=\beta b_1+(1-\beta)b_2$，则
\begin{align*}
u+v={}&\alpha\beta(a_1+b_1)+\alpha(1-\beta)(a_1+b_2)\\
&+(1-\alpha)\beta(a_2+b_1)+(1-\alpha)(1-\beta)(a_2+b_2).
\end{align*}
四个权重总和为1，组合的每个点都在 $A+B$ 中。这就是一般多重求和的缩小版；下面证明保留全部指标，且解释老师第58页的错误归一化为什么必须改。


\paragraph{原页55：凸包与和的交换及证明纠错。}凸包是所有包含 $S$ 的凸集的交，等价于有限凸组合全体。闭凸包是它的闭包。证明 $\operatorname{co}(A+B)=\operatorname{co}A+\operatorname{co}B$：一边把 $\sum\lambda_i(a_i+b_i)$ 分成两凸组合；另一边若 $a=\sum_i\alpha_i a_i,b=\sum_j\beta_j b_j$，则 $a+b=\sum_{i,j}\alpha_i\beta_j(a_i+b_j)$，积权重非负且总和1。有限多集合用权重 $\prod_k\lambda_{ki_k}$。第58页直接把总权重 $K$ 归一化会把目标点也除以 $K$，因此不成立；第59页积权重论证才正确。



\subsection{命题与证明：原页 57、58、59}

\textbf{命题。}$\operatorname{co}(\sum_i C_i)=\sum_i\operatorname{co}(C_i)$（有限和）。

\paragraph{老师的证明思路。}正向分配有限和；反向以乘积权重重组凸组合，不能简单归一化改变目标点。

\paragraph{逐步推导。}
记 $S=\sum_{k=1}^m S_k$。若有空因子，两边均空；以下均非空。凸包是有限凸组合的全体。
\paragraph{原页57：$\operatorname{co}S\subseteq\sum_k\operatorname{co}S_k$}
取 $x\in\operatorname{co}S$，则 $x=\sum_{i=1}^r\lambda_i z_i$，$z_i\in S$，$\lambda_i\ge0,\sum_i\lambda_i=1$。对每个 $z_i$ 选分解 $z_i=\sum_kz_{ki}$，$z_{ki}\in S_k$。有限和可交换次序：
\[x=\sum_i\lambda_i\sum_kz_{ki}=\sum_k\left(\sum_i\lambda_i z_{ki}\right).\]
每个括号中的点由 $S_k$ 内点按非负和1权重组合，属于 $\operatorname{co}S_k$，所以 $x$ 属于右侧和集。
\paragraph{原页58--59：反向必须使用乘积权重}
取 $x=\sum_k x_k$，每个 $x_k\in\operatorname{co}S_k$。分别表示
\[x_k=\sum_{i_k=1}^{r_k}\lambda_{k i_k}z_{k i_k},\quad \lambda_{k i_k}\ge0,\quad\sum_{i_k=1}^{r_k}\lambda_{k i_k}=1.\]
不能把所有权重总和 $m$ 直接归一化，那会把点 $x$ 也除以 $m$。对每个多重指标 $I=(i_1,\ldots,i_m)$，定义 $w_I=\prod_{k=1}^m\lambda_{k i_k}$，$z_I=\sum_{k=1}^m z_{k i_k}\in S$。权重总和为
\[\sum_{i_1=1}^{r_1}\cdots\sum_{i_m=1}^{r_m}w_I=\prod_{k=1}^m\left(\sum_{i_k=1}^{r_k}\lambda_{k i_k}\right)=1.\]
再验证组合的点\emph{恰为} $x$：
\begin{align*}
\sum_Iw_Iz_I
&=\sum_{k=1}^m\sum_{i_k=1}^{r_k}\lambda_{k i_k}z_{k i_k}
\prod_{j\ne k}\left(\sum_{i_j=1}^{r_j}\lambda_{j i_j}\right)\\
&=\sum_{k=1}^m\sum_{i_k=1}^{r_k}\lambda_{k i_k}z_{k i_k}
=\sum_kx_k=x.
\end{align*}
因此 $x$ 是 $S$ 点的合法凸组合，属 $\operatorname{co}S$，得到反向包含。两个方向合成等号。这个展开保留老师第59页所有多重求和的作用，同时说明为何那些额外乘上的和都等于1，不改变原点。


\clearpage\section{内积、超平面与 Minkowski 定理}\label{g:109:60}

\textbf{对应原课件：109，第 60--64 页。}


\paragraph{内积与超平面。}
$x\cdot y=\sum_{i=1}^N x_i y_i$ 是一个数，逐坐标乘法、有限求和连续，所以内积联合连续。$x\cdot y=0$ 称正交。
给 $p\ne0,c\in\R$，$H=\{z:p\cdot z=c\}$ 是法向量为 $p$ 的超平面。若 $x\in H$，则 $H=\{z:p\cdot(z-x)=0\}$。同时把 $p,c$ 乘同一个非零数不会改变平面。老师称超平面的交为 linear manifold；这里一般是仿射子空间，未必经过原点。

闭半空间 $\{z:p\cdot z\le c\}$ 闭：它是连续函数 $z\mapsto p\cdot z$ 对闭区间 $(-\infty,c]$ 的逆像。它凸：若 $p\cdot z_i\le c$，则 $p\cdot[(1-t)z_1+tz_2]\le(1-t)c+tc=c$。超平面是两个相反闭半空间的交，所以也闭凸；它们的任意交同样闭凸。
\paragraph{原页63--64的命题。}
超平面 $H$ bounding for $K$ 指 $K$ 在它的一侧闭半空间内，允许有点落在平面上。Minkowski 定理说：$K\subseteq\R^N$ 凸，给 $x\in\R^N$，存在经过 $x$ 且界住 $K$ 的超平面，当且仅当 $x\notin\Int K$。$x$ 可以在集合外；如果还在 $\overline K$，该平面是支撑平面。

原课件这里只陈述结论和图，没有展开证明。为了让这份讲义自足，后面的“补充证明”给完整有限维证明；它不冒充老师原有的证明路线。


\paragraph{原页60：内积、超平面与半空间。}$p\cdot x=\sum_i p_ix_i$ 为内积。固定非零 $p$，$H=\{x:p\cdot x=c\}$ 为超平面，$p$ 为法向量；非零条件避免整个空间或空集。两闭半空间由 $\ge c$ 与 $\le c$ 定义。线性函数连续，弱不等式在极限下保持，故闭；线性组合保持同侧不等式，故凸。




\subsection*{补充证明：原页63的 Minkowski 定理}
以下 $N\ge1$。要证的是存在 $p\ne0$，使 $p\cdot(z-x)\le0$ 对全部 $z\in K$ 成立；相应经过 $x$ 的平面是 $p\cdot(z-x)=0$。这把几何目标翻译成了一个可验证的不等式。

\paragraph{第一步：有平面则不能是内点。}
若 $x$ 是内点，就有球 $N(x,r)\subseteq K$。取 $z=x+(r/2)p/\|p\|$，它在这个球内，故在 $K$，但
\[p\cdot(z-x)=(r/2)\|p\|>0,
\]
与所有点都满足非正不等式矛盾。若原来在另一侧，把法向量变为 $-p$ 即可，因此必要性完成。

\paragraph{第二步：闭凸集外一点的最近点存在。}
暂设 $C$ 非空闭凸、$q\notin C$。取 $a\in C$，设 $R=\|q-a\|$。为了寻找最小距离，只需在 $C\cap\overline B(q,R)$ 中找：球外点比已知候选 $a$ 更远。这个交非空闭有界，所以紧；连续函数 $z\mapsto\|q-z\|$ 在其上达到最小值，记最近点为 $b$。因此它也是整个 $C$ 上的最近点。$q\notin C$ 保证 $q\ne b$。

任取 $z\in C$ 和 $0<t\le1$，凸性给 $b+t(z-b)\in C$，最近性给
\begin{align*}
\|q-b\|^2&\le\|q-b-t(z-b)\|^2\\
&=\|q-b\|^2-2t(q-b)\cdot(z-b)+t^2\|z-b\|^2.
\end{align*}
减去共同项并除以正数 $t$，得 $2(q-b)\cdot(z-b)\le t\|z-b\|^2$。让 $t\downarrow0$，得到
\[(q-b)\cdot(z-b)\le0\quad(z\in C).\tag{*}
\]
所以令 $p=q-b$，又有
\[p\cdot(z-q)=p\cdot(z-b)-\|p\|^2\le-\|p\|^2<0.
\]
这说明\emph{经过外点 $q$} 的平面也能界住 $C$。

\paragraph{第三步：闭凸集边界点的支撑平面。}
设 $x\in C\setminus\Int C$。它的每个邻域都含 $C$ 外点，可选 $q_n\notin C$ 且 $\|q_n-x\|<1/n$。为每个 $q_n$ 取最近点 $b_n\in C$。因 $x$ 本身是候选，
\[\|q_n-b_n\|\le\|q_n-x\|,
\quad\|b_n-x\|\le\|b_n-q_n\|+\|q_n-x\|<2/n.
\]
故 $b_n\to x$。把非零法向量归一化为 $p_n=(q_n-b_n)/\|q_n-b_n\|$，单位球面紧，抽子列 $p_{n_k}\to p$，且 $\|p\|=1$。将 $(*)$ 除以正的范数，得到 $p_{n_k}\cdot(z-b_{n_k})\le0$。对\emph{固定任意} $z\in C$ 取极限，内积连续给 $p\cdot(z-x)\le0$。$p\ne0$，即所需平面。这里不需要给每个 $z$ 重抽子列：子列早已只由 $p_n$ 选定。

\paragraph{第四步：原命题没要求 $K$ 闭，需要补这个环节。}
令 $C=\overline K$，它闭凸。我们需要保证 $x\notin\Int K$ 时也 $x\notin\Int C$。对凸集确有
\[\Int\overline K=\Int K.\tag{**}
\]
一边由 $K\subseteq\overline K$ 直接成立。证明另一边，取 $x\in\Int\overline K$，选球 $N(x,r)\subseteq\overline K$。取很小 $a>0$，使下面 $N+1$ 个点全在球中：
\[v_0=x-a\sum_{j=1}^Ne_j,\qquad v_i=x+ae_i\quad(1\le i\le N).
\]
这些点仿射独立，并且 $x=(v_0+\cdots+v_N)/(N+1)$，所有权重严格正。因此 $x$ 在该单纯形内部。每个 $v_i\in\overline K$，可选 $z_i\in K$ 任意接近 $v_i$；足够接近时仍仿射独立，且表示 $x$ 的重心坐标仍严格正。这个“仍”可由线性方程检查：把增广列 $(z_i,1)$ 组成方阵，其行列式在 $v_i$ 处非零，故小扰动下仍非零；解 $\sum_i\lambda_i z_i=x,\sum_i\lambda_i=1$ 由逆矩阵连续依赖各 $z_i$，因此各 $\lambda_i$ 仍接近正数 $1/(N+1)$。进一步把 $x$ 小幅移动时，权重也仍为正，所以 $x$ 的一个小球包含于 $\conv\{z_0,\ldots,z_N\}$。凸性使这个凸包包含于 $K$，故 $x\in\Int K$，证明 $(**)$。

现在 $x\notin\Int C$：若 $x\notin C$，用第二步；若 $x\in C$，用第三步。得到界住 $C$ 的平面，自然也界住 $K\subseteq C$。若 $K=\varnothing$，任意经过 $x$ 的超平面都符合条件，单独成立。充分性完成。

\paragraph{考试如何定位这份补充。}
课件原页只有定理陈述，以上最近点、单位法向量子列与凸集内部引理是自足补充。若老师只要求引用本定理，可直接写准确条件与结论；若要求证明，则不能把非闭凸集情况默默当作闭集。


\clearpage\section{极锥与坐标偏序}\label{g:109:65}

\textbf{对应原课件：109，第 65--66 页。}


\paragraph{极锥定义与完整验证。}
老师的符号 $C^\circ$ 在本页是\emph{极锥}，不是拓扑内部：
\[C^\circ=\{p:\forall x\in C,\ p\cdot x\le0\}
=\bigcap_{x\in C}\{p:p\cdot x\le0\}.
\]
每个因子是闭半空间（$x=0$ 时是整个空间），所以交闭且凸。$0\cdot x=0$ 给 $0\in C^\circ$；若 $p\in C^\circ,t\ge0$，则 $(tp)\cdot x=t(p\cdot x)\le0$，故是顶点0的锥。例如 $C=\R_+^N$ 的极锥是 $\R_-^N$：用各坐标单位向量测试得 $p_i\le0$，反向各项乘积非正即可。
\paragraph{坐标序。}
$x\le y$ 表示所有坐标 $x_i\le y_i$；$x<y$ 表示前一条件且至少一个严格；$x\ll y$ 表示每个坐标严格。比如 $(0,1)<(0,2)$ 但不满足 $\ll$；$(1,0)$ 与 $(0,1)$ 不可比较。这是偏序而不是全序。非负正交象限记为 $\Omega=\{x:x\ge0\}$，不是只有每个坐标严格正的点。


\clearpage\section{重新审视渐近方向与衰退方向}\label{g:109:67}

\textbf{对应原课件：109，第 67--75 页。}


\paragraph{第67--75页在重新组织同一套定义。}
截尾闭锥强调“任何远处阈值都留下的方向”；序列锥强调“存在一条点列及标度”；衰退锥强调“每个基点和每个非负步长”。原页72--74的几何说明：两条射线平移后重合，就称它们方向相同；两个非零向量方向相同当且仅当互为\emph{正}倍数，负倍数表示反方向。0没有几何朝向，但为了锥的代数性质仍加入方向向量集合。

这里不能把“无界集合”自动画成“已有射线”。一般无界集可以只有离散点或曲线；非空\textbf{闭凸}集的特殊性，才保证渐近方向能变成从每个基点出发的完整射线。后面的定理正是把这段直观变成证明。
\[
AC=\bigcap_{M\ge0}\Gamma(C^M),\qquad
C_\infty=\{\lim_k x_k/t_k:x_k\in C,\ t_k\to\infty\},
\]
\[0^+C=\{d:a+td\in C\text{ 对所有 }a\in C,t\ge0\}.
\]
第二个公式只收集确实存在的有限向量极限，不是说任意比值列都会收敛。


\paragraph{原页67：三种无穷远方向及量词。}截尾定义先保留 $\norm x\ge M$ 的点，再取所生成最小闭锥，最后对全部 $M$ 取交。序列定义为 $d\in C_\infty\iff\exists x_n\in C,\ t_n\to+\infty:x_n/t_n\to d$。衰退锥为 $\operatorname{rec}C=\{d:\forall x\in C,\forall t\ge0,x+td\in C\}$，要求所有起点都可以一直沿该方向走；只在非空闭凸情形才保证与渐近锥一致。单位方向归一化必须选范数趋于无穷的子列，而不仅是声称任意无界序列范数都趋无穷。截尾最小闭锥不是凸锥，闭包横线不能漏。



\subsection{命题与证明：原页 69、70}

\textbf{命题。}截尾集合及其生成闭锥随截尾阈值增大而递减；生成闭锥须保留闭包。

\paragraph{老师的证明思路。}先比较原截尾集包含，再用生成闭锥的单调性取交。

\paragraph{逐步推导。}
定义截尾集 $S^M=\{x\in S:\|x\|\ge M\}$。若 $M_1\le M_2$，任何满足更强条件 $\|x\|\ge M_2$ 的点也满足较弱条件，所以 $S^{M_2}\subseteq S^{M_1}$。
令 $\Gamma(E)$ 为包含 $E$ 且含0的最小闭锥。若 $E\subseteq F$，每个含 $F$ 的闭锥也含 $E$；尤其 $\Gamma(F)$ 是 $E$ 的一个闭锥超集，由最小性 $\Gamma(E)\subseteq\Gamma(F)$。应用于截尾集，得 $\Gamma(S^{M_2})\subseteq\Gamma(S^{M_1})$。所以这族闭锥随阈值递减，其集合列极限是交 $A S=\bigcap_M\Gamma(S^M)$。
老师再写生成锥的公式时必须保留闭包：
\[\Gamma(E)=\overline{\{\alpha x:\alpha\ge0,x\in E\}}\]
对非空 $E$ 成立。理由是未闭包的非负倍乘集已为锥且含 $E$；其闭包仍为锥，因为固定倍乘连续；任意闭锥超集又必须包含这些倍乘及其闭包，所以它确为最小。若 $E$ 空，则按“含0”的约定最小锥为 $\{0\}$，须单独处理。不能把生成闭锥误当凸包，也不能漏闭包。

原课件下一页还解释为何归一化能产生方向。若 $x_n\in C$ 且 $\|x_n\|\to\infty$，令 $u_n=x_n/\|x_n\|$。这些点全在有限维单位球面，紧性给子列 $u_{n_k}\to u$，而范数连续给 $\|u\|=1$，故极限非零。取 $t_k=\|x_{n_k}\|\to\infty$，有 $x_{n_k}/t_k\to u$，所以 $u\in C_\infty$。若原列仅无界而不保证范数趋无穷，须先逐层选 $\|x_n\|>n$ 的子列，再作这一构造。


\clearpage\section{渐近锥闭性、闭包不变与锥的特例}\label{g:109:76}

\textbf{对应原课件：109，第 76--80 页。}


\paragraph{这一组基本性质都不要求原集合凸。}
对非空 $C$，$C_\infty$ 是闭锥；$(\overline C)_\infty=C_\infty$；若 $C$ 本身是锥，则 $C_\infty=\overline C$，仅在 $C$ 已闭时才等于 $C$。
最费思考的是闭性：每个 $d_j$ 自己配有一条见证列，不能把所有见证的第 $j$ 项随手拼接。应从第 $j$ 条见证里选一个足够晚的项，同时保证标度大于 $j$、误差小于 $1/j$。下面把这个双指标选择完整写出。


\paragraph{原页76：渐近锥闭性及闭包不变。}非空 $C$ 时固定 $x\in C$ 并取 $t_n=n$ 得0属于渐近锥。正倍数通过把 $t_n$ 换为 $t_n/\lambda$ 得到。闭性用对角选取：若 $d_j\to d$，对每个 $j$ 从其见证序列选一点使 $t_j>j$ 且 $\norm{x_j/t_j-d_j}<1/j$；三角不等式给 $x_j/t_j\to d$。闭包不改变锥：用 $x_j'\in C$ 逼近 $x_j\in\overline C$ 至误差 $1/j$，除以 $t_j$ 后误差更趋0。若 $C$ 本身是闭锥，取 $x_n=nd,t_n=n$ 得 $C\subset C_\infty$，反向用每个 $x_n/t_n\in C$ 与闭性；非闭锥的结论为 $C_\infty=\overline C$。



\subsection{命题与证明：原页 77、78}

\textbf{命题。}非空集合的序列式渐近锥是闭锥。

\paragraph{老师的证明思路。}缩放构造锥性，选择每层精度与标度的对角点构造闭性。

\paragraph{逐步推导。}
设 $C\ne\varnothing$。\textbf{先验证锥性：}固定 $a\in C$，取 $x_n=a,t_n=n$，得到 $0\in C_\infty$。若 $d\in C_\infty$，有 $x_n/t_n\to d$。给任意 $\lambda>0$，新标度 $s_n=t_n/\lambda\to\infty$，同一原点列满足 $x_n/s_n=\lambda x_n/t_n\to\lambda d$，故正倍数仍在锥中；$\lambda=0$ 已单独处理。
\paragraph{原页78：闭性需要对角选取}
设 $d_j\in C_\infty,d_j\to d$。对每个 $j$，定义给一组自己的见证列 $x_{j,k}\in C,t_{j,k}\to\infty$ 且 $x_{j,k}/t_{j,k}\to d_j$。从第 $j$ 组中选一个足够晚的指标 $k(j)$，同时保证 $t_{j,k(j)}>j$ 和 $\|x_{j,k(j)}/t_{j,k(j)}-d_j\|<1/j$。令 $z_j=x_{j,k(j)}$、$s_j=t_{j,k(j)}$。则 $s_j>j\to\infty$，且
\[\|z_j/s_j-d\|\le\|z_j/s_j-d_j\|+\|d_j-d\|<1/j+\|d_j-d\|\to0.\]
所以 $z_j\in C,s_j\to\infty,z_j/s_j\to d$，恰为 $d\in C_\infty$ 的见证。每个锥内收敛列极限仍在锥内，故闭。直接把双指标都写成同一个 $n$ 而不选择 $k(j)$，不能保证同时满足精度与标度要求。


\subsection{命题与证明：原页 79}

\textbf{命题。}取闭包不改变渐近锥。

\paragraph{老师的证明思路。}用原集点逼近闭包中的见证点，控制除以标度后的误差。

\paragraph{逐步推导。}
因 $C\subseteq\overline C$，单调性先给 $C_\infty\subseteq(\overline C)_\infty$。反向取 $d\in(\overline C)_\infty$，有 $x_n\in\overline C,t_n\to\infty,x_n/t_n\to d$。闭包定义使每个 $n$ 可选 $z_n\in C$ 满足 $\|z_n-x_n\|<1/n$。注意选点随 $n$ 变化，且只要求指定误差内，不要求最近点存在。
充分大 $t_n>0$，误差估计给
\[\|z_n/t_n-d\|\le\|z_n-x_n\|/t_n+\|x_n/t_n-d\|<1/(nt_n)+\|x_n/t_n-d\|\to0.\]
因此原集合 $C$ 内的 $z_n$ 用同一标度也给方向 $d$，故 $d\in C_\infty$。得到反向包含，合为等号。这里不能把闭包点直接说成原集合点，补出的逼近及除标度误差正是两者之间的桥梁。


\subsection{命题与证明：原页 80}

\textbf{命题。}闭锥的渐近锥等于自身；一般锥的渐近锥等于其闭包。

\paragraph{老师的证明思路。}一边选 $x_n=nd$，另一边用锥在缩放下封闭及闭性。

\paragraph{逐步推导。}
先设 $C$ 为闭锥。若 $d\in C$，由锥性 $x_n=nd\in C$，取 $t_n=n$，则 $x_n/t_n=d$，所以 $d\in C_\infty$。反向若 $d\in C_\infty$，取见证 $x_n/t_n\to d$；充分大 $t_n>0$，锥性给 $x_n/t_n\in C$，闭性给极限 $d\in C$。两包含得 $C_\infty=C$。
对一般不闭锥，闭包仍是锥：若 $x_n\in C,x_n\to x$，对固定非负 $\lambda$ 有 $\lambda x_n\in C$ 且趋于 $\lambda x$，因此 $\lambda x\in\overline C$。前条已证闭包不改变渐近锥，所以
\[C_\infty=(\overline C)_\infty=\overline C.\]
原页写 $C_\infty=C$ 时必须看清右侧闭包横线；一般不闭锥只能得到闭包，不能漏掉它。


\clearpage\section{三个渐近描述为什么等价}\label{g:109:81}

\textbf{对应原课件：109，第 81--83 页。}


\paragraph{这两条等价定理连接前后定义。}
第82页证明 $AC=C_\infty$；第83页把非零方向写成单位极限方向的正倍数。令
\[C_N=\left\{u:\exists x_k\in C,\ \|x_k\|\to\infty,
\ x_k/\|x_k\|\to u\right\}.
\]
正确的统一写法为 $C_\infty=\{0\}\cup\{\lambda u:\lambda>0,u\in C_N\}$。当 $C$ 有界时 $C_N=\varnothing$，但渐近锥仍有0；不能忘记这一边界情况。截尾式中的 $\Gamma$ 有闭包，故属于它只保证\emph{可逼近}的倍数表示，不能无理由改成精确表示。


\paragraph{原页81：截尾与序列定义等价。}非零 $d=x_n/t_n$ 的极限使 $\norm{x_n}\to\infty$，所以对任意 $M$，最终原点均属截尾集，其缩放极限落入截尾闭锥，故属于交。反向对非零 $d$，每个整数 $M$ 选 $\alpha_M x_M$ 距 $d$ 小于 $1/M$，且 $\norm{x_M}\ge M$。由 $\alpha_M\le(\norm d+1/M)/M\to0$，取 $t_M=1/\alpha_M\to\infty$，得到见证序列。零向量对非空集合单独处理。归一化方向全体再取非负倍数得到相同锥；有界集方向集为空时需单独按约定加入0。



\subsection{命题与证明：原页 82}

\textbf{命题。}截尾生成闭锥定义与序列式渐近锥定义等价。

\paragraph{老师的证明思路。}非零方向的见证点最终越过每个截尾阈值；反向逐层选近似点并证明缩放系数趋零。

\paragraph{逐步推导。}
设 $C\ne\varnothing$，$AC=\bigcap_{M\ge0}\Gamma(C^M)$。最小闭锥约定都含0，所以0在每项交中；固定原点除以 $n$ 又使0在序列式锥中。因此零方向单独完成，以下均取 $d\ne0$。
\paragraph{序列方向推出截尾方向}
若 $x_n/t_n\to d$，则范数趋于 $\|d\|>0$，充分大有 $\|x_n\|/t_n>\|d\|/2$，所以 $\|x_n\|\to\infty$。固定任意 $M$，最终 $x_n\in C^M$，且 $t_n>0$，因此 $x_n/t_n$ 在 $C^M$ 的非负倍乘集中，亦在其生成闭锥。这个闭锥含尾部及其极限，故 $d\in\Gamma(C^M)$。$M$ 任意，$d\in AC$。
\paragraph{截尾方向推出序列方向}
给 $d\in AC\setminus\{0\}$。对充分大的整数 $M$，因 $d$ 属于截尾\emph{闭}锥，选 $x_M\in C$，$\|x_M\|\ge M$，以及 $\alpha_M\ge0$，使 $\|d-\alpha_Mx_M\|<1/M$。当 $1/M<\|d\|/2$ 时，$\alpha_M$ 不可能0，否则误差等于 $\|d\|$。由三角不等式，
\[0<\alpha_M\le\frac{\|d\|+\|d-\alpha_Mx_M\|}{\|x_M\|}<\frac{\|d\|+1/M}{M}\to0.\]
令 $t_M=1/\alpha_M\to\infty$，则 $x_M/t_M=\alpha_Mx_M\to d$，且 $x_M\in C$，所以 $d$ 属于序列锥。两向均建立后才可在后续任意替换两种定义。
原页将属于闭包直接写成属于未闭倍乘集的地方应理解为此处的 $1/M$ 逼近，不能要求恰好表示。


\subsection{命题与证明：原页 83}

\textbf{命题。}无界集合的渐近方向可由无穷远单位方向及非负缩放描述；有界情形另补零方向约定。

\paragraph{老师的证明思路。}分离长度比例与单位方向，分别抽子列与取极限。

\paragraph{逐步推导。}
令 $C_N$ 为从 $C$ 中范数趋无穷点列归一化后可得到的单位极限集合。先取 $u\in C_N$，则存在 $x_n\in C,\|x_n\|\to\infty$，且 $x_n/\|x_n\|\to u$。以 $t_n=\|x_n\|$ 作为标度，定义直接给 $u\in C_\infty$；渐近锥对非负倍数封闭，所以任意 $\lambda u$ 也在其中。
反向取非零 $d\in C_\infty$，见证 $x_n/t_n\to d$。范数连续给 $r_n=\|x_n\|/t_n\to\|d\|>0$，所以 $\|x_n\|=r_nt_n\to\infty$。又
\[\frac{x_n}{\|x_n\|}=\frac{x_n/t_n}{\|x_n\|/t_n}\longrightarrow\frac d{\|d\|}=:u.\]
因此 $u\in C_N$ 且 $d=\|d\|u$，得非零方向的反向包含。原页再抽长度比例子列也可；它实际上已有正极限 $\|d\|$，故不能只写 $\lambda\ge0$ 后未经说明就除以它。
零方向需注意约定：若 $C$ 无界，则 $C_N$ 非空，取任何 $u$ 的0倍即得0；若 $C$ 非空有界，$C_N=\varnothing$，普通集合 $\{\lambda u:u\in C_N\}$ 为空，不能等于 $C_\infty=\{0\}$。需将生成锥约定额外含0，或把此等式限定于无界情形，再单独处理有界情形。


\clearpage\section{用序列重证有界性与平移性质}\label{g:109:84}

\textbf{对应原课件：109，第 84--87 页。}


\paragraph{原课件在这里用序列定义复习性质。}
这一组重新证明 $C$ 有界当且仅当 $C_\infty=\{0\}$、$(C+a)_\infty=C_\infty$，以及非空 $E$ 时 $C_\infty\subseteq(C+E)_\infty$。同样的定理在截尾语言里用包含关系，在序列语言里用见证序列转换。为方便按老师次序复习，本节再次给出完整证明，不要求翻回前半部分。


\subsection{命题与证明：原页 84}

\textbf{命题。}非空集合有界当且仅当其序列式渐近锥只有零向量。

\paragraph{老师的证明思路。}有界时除以趋无穷标度；无界时选范数趋无穷点列并单位化抽子列。

\paragraph{逐步推导。}
本条 $C\ne\varnothing$。若 $C$ 有界，取 $M$ 使任意 $x\in C$ 有 $\|x\|\le M$。任取 $d\in C_\infty$，定义给 $x_n\in C,t_n\to+\infty$ 且 $x_n/t_n\to d$。充分大 $t_n>0$，有 $\|x_n/t_n\|\le M/t_n\to0$，所以 $d=0$。反向0确实在锥中：固定 $x\in C$，取 $x_n=x,t_n=n$，则 $x/n\to0$。故锥恰为 $\{0\}$。
反向假设锥只有0，若 $C$ 无界，可对每个 $n$ 选 $x_n\in C$ 使 $\|x_n\|>n$。令 $u_n=x_n/\|x_n\|$，范数恒1。有限维单位球面紧，存在子列 $u_{n_k}\to u$，范数连续给 $\|u\|=1$，所以 $u\ne0$。
取对应 $t_k=\|x_{n_k}\|\to\infty$，则 $x_{n_k}/t_k=u_{n_k}\to u$，定义使 $u\in C_\infty$，与只有0矛盾。因此 $C$ 有界。不能说任意无界列范数都趋无穷，第一步的逐层选点专门保证这一点。


\subsection{命题与证明：原页 85}

\textbf{命题。}平移不改变渐近锥：$(C+a)_\infty=C_\infty$。

\paragraph{老师的证明思路。}把见证序列加减固定向量；除以趋无穷标度后平移项趋零。

\paragraph{逐步推导。}
任取 $d\in(C+a)_\infty$。按定义有 $z_n\in C+a,t_n\to\infty$ 且 $z_n/t_n\to d$。逐项写 $z_n=x_n+a$，其中 $x_n\in C$。则
\[\left\|\frac{x_n}{t_n}-d\right\|\le\left\|\frac{z_n}{t_n}-d\right\|+\frac{\|a\|}{t_n}\to0.\]
两项趋零的理由分别是见证序列的已知收敛、固定平移除以趋无穷标度。因此 $d\in C_\infty$，得到 $(C+a)_\infty\subseteq C_\infty$。
反向对任意 $d\in C_\infty$，取原见证 $x_n/t_n\to d$，令 $z_n=x_n+a\in C+a$。同一误差估计给 $z_n/t_n\to d$，故 $d\in(C+a)_\infty$。或者按老师写法，将第一包含中的集合换为 $C+a$、平移换为 $-a$，也得到反向。两包含合成等号。
这说明固定原点位置的改变被无穷尺度消掉，不是说原集合本身平移后相等。


\subsection{命题与证明：原页 86}

\textbf{命题。}平移不改变渐近锥：$(C+a)_\infty=C_\infty$。

\paragraph{老师的证明思路。}把见证序列加减固定向量；除以趋无穷标度后平移项趋零。

\paragraph{逐步推导。}
老师的等价链必须展开为见证序列的双向转换：
\begin{align*}
d\in C_\infty
&\iff\exists x_n\in C,\ t_n\to\infty:\ x_n/t_n\to d\\
&\iff\exists z_n=x_n+a\in C+a,\ t_n\to\infty:\ z_n/t_n\to d\\
&\iff d\in(C+a)_\infty.
\end{align*}
中间等价不是仅靠重命名，因为表达式多出了 $a/t_n$。正向若 $x_n/t_n\to d$，则 $z_n/t_n=x_n/t_n+a/t_n\to d+0=d$；反向若 $z_n/t_n\to d$，则 $x_n=z_n-a\in C$ 且 $x_n/t_n=z_n/t_n-a/t_n\to d$。$a$ 固定保证 $a/t_n\to0$，标度保持不变，所以两方向见证都合法。
最后截尾定义与序列定义已经证明等价时，可将 $C_\infty$ 换写为 $AC$，得到原页的记号。不能只对其中一个包含验证见证转换便称为等价链。


\subsection{命题与证明：原页 87}

\textbf{命题。}非空 $E$ 下，$C_\infty\subseteq(C+E)_\infty$。

\paragraph{老师的证明思路。}固定 $a\in E$，结合 $C+a\subseteq C+E$ 与平移不变性。

\paragraph{逐步推导。}
设 $E\ne\varnothing$，固定一点 $a\in E$。任意 $c\in C$ 都有 $c+a\in C+E$，故 $C+a\subseteq C+E$。渐近锥单调性给 $(C+a)_\infty\subseteq(C+E)_\infty$。平移不变性又给 $(C+a)_\infty=C_\infty$，于是所需包含成立。
如直接验证见证：给 $d\in C_\infty$ 的 $c_n/t_n\to d$，点 $c_n+a\in C+E$ 且 $(c_n+a)/t_n=c_n/t_n+a/t_n\to d$。这解释原页两句引用实际如何作用，不需要令所选 $a$ 随 $n$ 改变，也不需要 $E$ 有界。


\clearpage\section{乘积与交：不能交换的运算}\label{g:109:88}

\textbf{对应原课件：109，第 88--91 页。}


\paragraph{第89页等号错误的最短检验。}
在 $\R^2$ 中 $(M,0)$ 的总范数至少 $M$，却不满足第二坐标绝对值至少 $M$。所以“积集合的截尾”不等于“每个因子的截尾之积”。老师要证明的锥包含是真的，但该等号不能用于证明；应回到同一标度下的坐标投影，下面保留修正过程。
交集有界的证明也不能只从无限交没有方向就断言某一截尾为空，要补上单位球面紧性或已经证明的有界性判据。


\subsection{命题与证明：原页 88}

\textbf{命题。}有限积的渐近锥包含于各因子渐近锥的积。

\paragraph{老师的证明思路。}对共同标度下的向量极限逐坐标取极限；原页截尾等式的错误另外说明。

\paragraph{逐步推导。}
设 $C=\prod_{i=1}^mC_i$。若积为空，包含按相应空锥约定单独成立；以下取 $d\in C_\infty$，由定义选 $x_n=(x_{1n},\ldots,x_{mn})\in C$ 与同一 $t_n\to\infty$，满足 $x_n/t_n\to d=(d_1,\ldots,d_m)$。
对每个固定 $i$，积归属给 $x_{in}\in C_i$，坐标投影连续给 $x_{in}/t_n\to d_i$。这正是 $d_i\in(C_i)_\infty$ 的完整见证，故全部坐标合起来得到
\[d\in\prod_i(C_i)_\infty.\]
由任意 $d$ 得 $C_\infty\subseteq\prod_i(C_i)_\infty$。0方向也被此推理覆盖，无须跳过。
反向不能从各自见证直接拼接，因为不同因子的 $t_n$ 可能不相同；本条只证明包含。原页89另用截尾的证明存在错误，见其单列辨错步骤。


\subsection{命题与证明：原页 89}

\textbf{命题。}有限积的渐近锥包含于各因子渐近锥的积。

\paragraph{老师的证明思路。}对共同标度下的向量极限逐坐标取极限；原页截尾等式的错误另外说明。

\paragraph{逐步推导。}
原页欲证积的渐近锥包含于各因子锥的积，但其截尾推导的第一条关键等式错误：
\[\{(x_1,\ldots,x_m)\in\prod_iC_i:\|(x_1,\ldots,x_m)\|\ge M\}\ne\prod_i\{x_i\in C_i:\|x_i\|\ge M\}\]
一般不相等。例如 $(M,0)$ 的总范数达到 $M$，第二坐标却未达阈值。随后将“共同倍数缩放的积”当成“各自独立倍数缩放的积”的等号也未成立，共同标度不能凭空分开。
保留老师想要的逐分量方向结论，可在序列定义下完整修正：取 $d\in(\prod_iC_i)_\infty$，选同一标度 $t_n\to\infty$ 与 $x_n=(x_{in})\in\prod_iC_i$，使 $x_n/t_n\to d=(d_i)$。对固定 $i$，$x_{in}\in C_i$ 且投影收敛 $x_{in}/t_n\to d_i$，因此 $d_i\in(C_i)_\infty$。对全部 $i$ 成立，故 $d\in\prod_i(C_i)_\infty$。
这与前页88的合法原路线相同，是补正错误所必需的替代；不能继续沿不正确截尾等号推到结论后称其为证明。


\subsection{命题与证明：原页 90、91}

\textbf{命题。}若各集合渐近锥的交至多为零方向，则原集合的交有界。

\paragraph{老师的证明思路。}交集包含于每一因子，先用锥单调性，再用有界性判据。

\paragraph{逐步推导。}
设 $C=\bigcap_\gamma C_\gamma$，且 $\bigcap_\gamma(C_\gamma)_\infty=\{0\}$。若 $C$ 为空，空集本来有界；以下非空。
固定任意 $\gamma$，有 $C\subseteq C_\gamma$。渐近锥单调性给 $C_\infty\subseteq(C_\gamma)_\infty$。对所有 $\gamma$ 都成立，故
\[C_\infty\subseteq\bigcap_\gamma(C_\gamma)_\infty=\{0\}.\]
因为 $C$ 非空，固定一点并除以 $n$ 可证 $0\in C_\infty$。因此不是仅有包含，而是 $C_\infty=\{0\}$。由前条“非空集合有界当且仅渐近锥只有0”，得到 $C$ 有界。
如果按老师另一版写成截尾，零以外没有渐近方向便意味着不可能存在任意远的交集点；否则将其归一化，单位球面紧性会给一个非零方向，落在上述所有锥中，矛盾。不能直接从交锥只有0无理由写“某阈值后截尾为空”，中间正是这个有界性判据。


\clearpage\section{再走一遍线性像定理}\label{g:109:92}

\textbf{对应原课件：109，第 92--96 页。}


\paragraph{第二遍读时把四个箭头自己补上。}
像有界 $\Rightarrow$ 原像有界：反证、归一化、核条件。原像有界 $+$ 集合闭 $\Rightarrow$ 像闭：抽子列、取极限。原集合方向 $\Rightarrow$ 像方向：对见证用线性映射。像方向 $\Rightarrow$ 原集合方向：先证缩放原像有界，再抽子列。下面保留完整四段，原页93--96分别对应这四步。


\paragraph{原页92：线性像定理再证明。}本段重复第37--41页的重要定理。核心核条件为 $\ker L\cap C_\infty=\{0\}$。若 $Lx_n$ 有界而 $x_n$ 无界，取范数趋无穷子列并单位化，极限是非零渐近方向且被 $L$ 送到0，矛盾。因此像序列收敛时原像可抽收敛子列，闭 $C$ 与连续 $L$ 给像闭。锥像的容易包含是对 $x_n/t_n$ 直接取 $L$；困难包含须用第41页的缩放反证，不可仅引用原集合上的有界性。



\subsection{命题与证明：原页 93、94、95、96}

\textbf{命题。}闭 $C$ 满足 $\ker L\cap C_\infty=\{0\}$ 时，线性像 $L(C)$ 闭且 $L(C_\infty)=(L(C))_\infty$。

\paragraph{老师的证明思路。}归一化反证原像有界；抽子列证像闭；两向分别构造渐近方向。

\paragraph{逐步推导。}
设闭 $C\subseteq\R^N$，线性 $L:\R^N\to\R^K$ 满足 $\ker L\cap C_\infty=\{0\}$。有限维线性映射连续；例如矩阵系数的有限和给 $\|Lz\|\le M\|z\|$。
\paragraph{原页38/93：像有界迫使原像有界}
若有 $x_n\in C$ 的线性像有界，而原像列无界，取子列（重记）使 $r_n=\|x_n\|\to\infty$。单位向量 $u_n=x_n/r_n$ 再抽子列趋于单位 $u$。由渐近锥定义 $u\in C_\infty$。又
\[Lu=\lim_n L(x_n/r_n)=\lim_n(Lx_n)/r_n=0,\]
因为分子有界、分母趋无穷。故 $u\in\ker L\cap C_\infty$ 却 $\|u\|=1$，矛盾。得到所需有界性。
\paragraph{原页39/94：像闭}
取任意 $y_n\in L(C)$ 且 $y_n\to y$。每项选 $x_n\in C$ 使 $Lx_n=y_n$。收敛像列有界，上段使 $x_n$ 有界，故抽 $x_{n_k}\to x$。$C$ 闭给 $x\in C$，$L$ 连续给 $Lx=\lim_kLx_{n_k}=\lim_ky_{n_k}=y$。于是 $y\in L(C)$，像闭。原像不要求唯一，任意分解都受同一核条件约束。
\paragraph{原页40/95：容易方向 $L(C_\infty)\subseteq(L(C))_\infty$}
给 $d\in C_\infty$，取 $x_n\in C,t_n\to\infty,x_n/t_n\to d$。令 $y_n=Lx_n\in L(C)$，则 $y_n/t_n=L(x_n/t_n)\to Ld$，所以 $Ld$ 是像集的渐近方向。此方向不需要闭性及核条件。
\paragraph{原页41/96：反向，补上不能直接套用的缩放有界性}
任取 $y\in(L(C))_\infty$，选 $y_n\in L(C),t_n\to\infty,y_n/t_n\to y$，再选 $x_n\in C$ 使 $Lx_n=y_n$。需证 $x_n/t_n$ 有界；因为这些缩放点\emph{未必在 $C$ 中}，不能直接应用第一段。
反设缩放列无界，取子列使 $q_n=\|x_n\|/t_n\to\infty$，并进一步使 $u_n=x_n/\|x_n\|\to u$，$\|u\|=1$。这时 $\|x_n\|=q_nt_n\to\infty$，所以 $u\in C_\infty$，且
\[Lu=\lim_n\frac{Lx_n}{\|x_n\|}=\lim_n\frac{y_n/t_n}{q_n}=0.\]
又与核条件矛盾。故 $x_n/t_n$ 有界，可抽子列趋于 $d$。同一原点及标度满足定义，$d\in C_\infty$；连续性给 $Ld=\lim_nL(x_n/t_n)=y$。所以 $y\in L(C_\infty)$。两方向合成等号。原页略过的缩放有界证明是本定理最关键的补步。


\clearpage\section{闭和定理：线性映射法与直接子列法}\label{g:109:97}

\textbf{对应原课件：109，第 97--103 页。}


\paragraph{这次老师明确给两条路线。}
第98--100页先用积集与求和线性映射，验证核条件后调用线性像定理。第101--103页直接从收敛和列出发，假设一个分量无界，将分解全部除以同一大标度，得到相反非零方向，从而矛盾。第二种是第一种机制在两集合情形下的直接展开，不能省掉正半独立的使用位置。


\paragraph{原页97：闭和定理的两个证明。}积集加求和映射的证明已在第42页完整给出。两集合的直接证明：若 $z_n=a_n+b_n\to z$ 而 $a_n$ 无界，选 $r_n=\norm{a_n}\to\infty$ 及 $a_n/r_n\to d$，$\norm d=1$。由 $z_n/r_n\to0$ 得 $b_n/r_n\to-d$，于是 $d\in A_\infty,-d\in B_\infty$，违反正半独立。故 $a_n$ 有界，抽子列趋于 $a\in A$，再由 $b_n=z_n-a_n\to z-a\in B$，得 $z\in A+B$。有限多集合用积方法最直接，避免在归纳时漏检中间和的渐近锥条件。



\subsection{命题与证明：原页 98、99、100}

\textbf{命题。}闭集族渐近锥正半独立时，有限集合和闭，且和的渐近锥包含于各锥之和。

\paragraph{老师的证明思路。}将和表示为积集在线性求和映射下的像，检查其核与渐近锥交。

\paragraph{逐步推导。}
假设各闭 $C_i\subseteq\R^N$ 的渐近锥正半独立：若 $d_i\in(C_i)_\infty$ 且 $\sum_i d_i=0$，则每个 $d_i=0$。若某因子空，和为空，闭性直接成立；以下非空。
令 $C=\prod_{i=1}^mC_i$，并定义线性映射 $L(x_1,\ldots,x_m)=\sum_i x_i$。有限闭积闭，且 $L(C)=\sum_iC_i$。积锥包含已给
\[C_\infty\subseteq\prod_i(C_i)_\infty.\]
若 $d=(d_1,\ldots,d_m)\in\ker L\cap C_\infty$，则每个 $d_i\in(C_i)_\infty$，同时 $\sum_id_i=Ld=0$。正半独立性强迫全部 $d_i=0$，故 $d=0$。反向0确实在该交，所以核条件正好成立。应用线性像定理，$L(C)=\sum_iC_i$ 闭。
再证明老师随后的锥包含：线性像定理给 $(L(C))_\infty=L(C_\infty)$；积锥包含经 $L$ 的单调取像给
\[L(C_\infty)\subseteq L\left(\prod_i(C_i)_\infty\right)=\sum_i(C_i)_\infty.\]
所以 $(\sum_iC_i)_\infty\subseteq\sum_i(C_i)_\infty$。最后一个等号只是求和映射的像定义，不要求在一般情形积锥等于锥的积。核条件须先验证再用定理，不能仅说 $L$ 连续便断言闭像。


\subsection{命题与证明：原页 101、102、103}

\textbf{命题。}闭和定理的直接归一化反证版本。

\paragraph{老师的证明思路。}若分解中的一个分量无界，除以其范数后得到两个抵消的非零方向。

\paragraph{逐步推导。}
先看两个闭集 $A,B$。给任意 $z_n\in A+B,z_n\to z$，逐项选 $a_n\in A,b_n\in B$ 使 $z_n=a_n+b_n$。正半独立意味着不存在非零 $d\in A_\infty$ 同时 $-d\in B_\infty$。
若 $a_n$ 无界，先取子列使 $r_n=\|a_n\|\to\infty$，再从单位向量中抽子列使 $a_n/r_n\to d$，$\|d\|=1$。由于 $z_n$ 收敛有界，$z_n/r_n\to0$。于是
\[b_n/r_n=z_n/r_n-a_n/r_n\to-d.\]
同一标度 $r_n\to\infty$ 给 $d\in A_\infty,-d\in B_\infty$，它们相加为0却不为零，违反正半独立。因此 $a_n$ 有界。抽子列 $a_{n_k}\to a$，$A$ 闭使 $a\in A$；对应 $b_{n_k}=z_{n_k}-a_{n_k}\to z-a$，$B$ 闭使 $z-a\in B$。所以 $z=a+(z-a)\in A+B$，和闭。
老师说一般情形归纳，还须检查中间和的条件。若已知前 $m-1$ 个和闭且其渐近锥包含于前 $m-1$ 个锥之和，则中间和方向 $d$ 与最后方向 $e$ 抵消时，把 $d$ 分解为各旧锥方向，整体正半独立使所有分量及 $e$ 为0。因此二集条件可用。中间和的渐近锥包含可由前面的积集加求和线性映射定理建立。这样归纳才不是未检查条件的一句话。


\clearpage\section{渐近正则：把不同尺度统一}\label{g:109:104}

\textbf{对应原课件：109，第 104--108 页。}


\paragraph{新增定义只有一个量词变了。}
普通 $C_\infty$：\emph{存在} $t_k\to\infty$，并存在点列满足 $x_k/t_k\to d$。
强方向集合
\[C^1_\infty=\{d:\forall t_k\to\infty\ \exists x_k\in C,
\ x_k/t_k\to d\}
\]
要求对预先任意给定的标度都能选点。显然 $C^1_\infty\subseteq C_\infty$，因为可以特取 $t_k=k$。两者相等时称 $C$ 渐近正则。

老师先证非空凸集渐近正则，再证其渐近锥凸：两个方向原来可能使用不同标度，有了正则性就都能改用 $k$，然后把两条点列做同一权重的凸组合。\textbf{这里凸集无需闭；下一组与衰退锥相等才需要闭性。}


\paragraph{原页104：渐近正则性与凸渐近锥。}正则性把“存在缩放速度”增强为“每一种趋无穷速度都可匹配”。设 $x_j/s_j\to d$，固定 $a\in C$，给任意 $t_n\to\infty$ 选 $j(n)\to\infty$ 且 $s_{j(n)}\ge t_n$。由凸性 $z_n=a+(t_n/s_{j(n)})(x_{j(n)}-a)\in C$，且 $z_n/t_n=a/t_n+(x_{j(n)}-a)/s_{j(n)}\to d$。因此任意非空凸集渐近正则。对两方向使用共同 $t_n=n$ 的见证点，再作相同权重凸组合，就证明渐近锥凸；其闭锥性质已由第76页给出。



\subsection{命题与证明：原页 106、107}

\textbf{命题。}非空凸集渐近正则：任意趋无穷标度均能构造给定方向的见证点。

\paragraph{老师的证明思路。}从已有标度选更大的见证点，再沿固定基点的线段按比例缩回所需标度。

\paragraph{逐步推导。}
定义强方向集合 $C^1_\infty$：给任意标度 $t_n\to\infty$，都能找到 $z_n\in C$ 使 $z_n/t_n\to d$。显然强集合包含于普通渐近锥，因为可以特取 $t_n=n$。
反向设非空凸 $C$，给 $d\in C_\infty$，选一组旧见证 $x_k/s_k\to d$，$s_k\to\infty$。固定 $a\in C$，令 $u_k=(x_k-a)/s_k$，则 $u_k\to d$，且 $a+s_ku_k=x_k\in C$。
现在给\emph{任意另一条}标度 $t_n\to\infty$。对每个 $n$，选 $k(n)\ge n$ 且 $s_{k(n)}\ge t_n$；这由 $s_k\to\infty$ 保证。令
\[z_n=a+t_nu_{k(n)}=\left(1-\frac{t_n}{s_{k(n)}}\right)a+\frac{t_n}{s_{k(n)}}x_{k(n)}.\]
比例在 $[0,1]$，两端点在 $C$，凸性给 $z_n\in C$。又 $k(n)\to\infty$，所以
\[z_n/t_n=a/t_n+u_{k(n)}\to0+d=d.\]
任意标度均构造出见证，故 $d\in C^1_\infty$，两方向相等。选 $k(n)\ge n$ 的条件保证取到的方向近似确实越来越精确，而不能仅选任意满足大标度的一项后固定不动。


\subsection{命题与证明：原页 108}

\textbf{命题。}非空凸集的渐近锥是闭凸锥。

\paragraph{老师的证明思路。}利用渐近正则性统一两组见证点的标度，再作凸组合。

\paragraph{逐步推导。}
已由前面的对角论证知道 $C_\infty$ 闭且为锥；本条只需补凸性，但不能据原集合凸就直接宣布极限方向也凸。取任意 $d_1,d_2\in C_\infty$ 及 $\lambda\in[0,1]$。
前条非空凸集渐近正则允许我们对\emph{共同}标度 $t_n=n$ 分别选 $x_{1n},x_{2n}\in C$，使 $x_{in}/n\to d_i$。若只用普通定义，各方向本来的标度未必相同，这正是先证正则性的目的。
令 $z_n=\lambda x_{1n}+(1-\lambda)x_{2n}$。凸性给 $z_n\in C$，且有限线性运算保极限：
\[z_n/n=\lambda(x_{1n}/n)+(1-\lambda)(x_{2n}/n)\to\lambda d_1+(1-\lambda)d_2.\]
用序列锥定义得到这个凸组合也在 $C_\infty$。因端点与权重任意，锥凸。闭、锥、凸分别来自不同论证，不能把它们视为同一性质。


\clearpage\section{闭凸集的渐近锥等于衰退锥}\label{g:109:109}

\textbf{对应原课件：109，第 109--113 页。}


\paragraph{核心定理与不能删掉的条件。}
任意非空集合都有 $0^+C\subseteq C_\infty$。非空闭凸集则 $0^+C=C_\infty$。
第109页最后一句若删掉 closed 仍宣称等号，就是错误；第112页正是用来说明闭性必要的例子。

例子为 $C=\{(u,v):0\le u<1,v\ge0\}\cup\{(1,0)\}$。它凸：两个左条带点的组合仍在条带；若一个端点为 $(1,0)$，除该端点自身外的组合第一坐标严格小于1，第二坐标非负。它不闭，因为 $(1-1/n,1)\to(1,1)$ 不在集合。
正确答案是 $0^+C=\{0\}$，$C_\infty=\{(0,t):t\ge0\}$。不是空衰退锥，也不是整个竖直直线；0总满足衰退定义，且原点的第二坐标非负使渐近方向第二坐标也非负。后面的详细证明分别检查每个分量。


\paragraph{原页109：衰退锥相等与非闭反例。}任意非空集合有 $\operatorname{rec}C\subset C_\infty$：固定 $a\in C$，用 $x_n=a+nd$。反向需要闭凸：若 $x_n/t_n\to d$，对任意 $a\in C$ 和固定 $t>0$，充分大时 $a+(t/t_n)(x_n-a)\in C$，取极限得 $a+td\in C$。所以 $\operatorname{rec}C=C_\infty$。第112页反例 $C=([0,1)\times[0,\infty))\cup\{(1,0)\}$ 中，衰退方向必须让 $(1,0)$ 及任意 $(0,y)$ 都留在集合，迫使两个方向分量均0，故衰退锥是 $\{0\}$，不是空集；渐近锥是 $\{(0,v):v\ge0\}$，不是整条竖直线。闭凸时 $C+C_\infty=C$ 由衰退定义及0在锥中立即得到。



\subsection{命题与证明：原页 110、111}

\textbf{命题。}衰退锥包含于渐近锥；非空闭凸时两锥相等。

\paragraph{老师的证明思路。}一边沿射线取点；另一边以凸组合逼近任意基点加任意长度方向，再用闭性。

\paragraph{逐步推导。}
定义 $\operatorname{rec}C=\{d:\forall a\in C,\forall t\ge0,a+td\in C\}$。先仅设 $C$ 非空。
\paragraph{衰退方向总是渐近方向}
若 $d\in\operatorname{rec}C$，固定 $a\in C$，令 $x_n=a+nd\in C$，$t_n=n$。则 $x_n/t_n=a/n+d\to d$，故 $d\in C_\infty$。
\paragraph{闭凸时反向成立}
设 $C$ 另闭凸，$d\in C_\infty$，选 $x_n\in C,s_n\to\infty,x_n/s_n\to d$。要证衰退条件，必须\emph{任取}基点 $a\in C$ 及固定长度 $t>0$。充分大 $s_n\ge t$，比例 $t/s_n\in[0,1]$，所以
\[z_n=\left(1-\frac t{s_n}\right)a+\frac t{s_n}x_n\in C.\]
而 $z_n=a+t(x_n/s_n)-ta/s_n\to a+td$。闭性给 $a+td\in C$。$t=0$ 时是原点 $a$，自动成立；基点与长度任意，所以 $d\in\operatorname{rec}C$。
两包含合成 $\operatorname{rec}C=C_\infty$。若只假设凸而不闭，最后从极限回到集合内这一步不能做；原页末句此时只能对 $\overline C$ 说 $\operatorname{rec}(\overline C)=(\overline C)_\infty=C_\infty$，不能去掉闭包横线。


\subsection{命题与证明：原页 112}

\textbf{命题。}非闭凸集的衰退锥与渐近锥可以不同：课件条带加端点例证。

\paragraph{老师的证明思路。}逐坐标检查哪些射线能从所有基点出发，再直接构造无穷远缩放序列。

\paragraph{逐步推导。}
令 $C=([0,1)\times[0,\infty))\cup\{(1,0)\}$。先算衰退锥。若 $d=(d_1,d_2)$ 是衰退方向，从任意基点沿它走所有长度都须保持第一坐标在 $[0,1]$。若 $d_1>0$，从 $(0,0)$ 走足够远会超过1；若 $d_1<0$，从 $(0,0)$ 立即变负。因此 $d_1=0$。
若 $d_2<0$，从 $(0,0)$ 走后第二坐标负；若 $d_2>0$，从 $(1,0)$ 走到 $(1,td_2)$，该点第一坐标恰1而第二坐标正，不在 $C$。故 $d_2=0$，于是 $\operatorname{rec}C=\{0\}$，不是原页写的空集，因为0方向永远合法。
再算渐近锥：若 $(u_n,v_n)\in C,t_n\to\infty$，第一坐标 $0\le u_n\le1$，所以任何缩放极限第一坐标都0；第二坐标 $v_n\ge0$，极限不可能负，故 $C_\infty\subseteq\{(0,b):b\ge0\}$。反向给任意 $b\ge0$，取 $(u_n,v_n)=(0,nb)\in C$、$t_n=n$，缩放恒为 $(0,b)$，故该射线全在渐近锥中。
所以渐近锥为竖直非负射线，而非原页写的整条直线。这严格显示闭性删除后两个锥确可不同。


\subsection{命题与证明：原页 113}

\textbf{命题。}非空闭凸集满足 $C+C_\infty=C$。

\paragraph{老师的证明思路。}0方向给一个包含；反向用渐近见证点的凸组合趋于所需和点。

\paragraph{逐步推导。}
设非空闭凸 $C$。0在渐近锥中，所以对任意 $c\in C$，$c=c+0\in C+C_\infty$，得到 $C\subseteq C+C_\infty$。
反向任取 $x\in C+C_\infty$，写 $x=c+d$，$c\in C,d\in C_\infty$。取方向见证 $z_n\in C,t_n\to\infty,z_n/t_n\to d$。充分大 $t_n\ge1$，故
\[w_n=(1-1/t_n)c+(1/t_n)z_n\in C\]
为合法凸组合。且 $w_n=c-c/t_n+z_n/t_n\to c+d=x$。闭性给 $x\in C$。两包含证等号。
这保留老师的直接凸组合路线，即使已经证明衰退锥相等也不将本条压缩成“由定义显然”。凸性用来给每个 $w_n$ 归属，闭性用来给最终 $x$ 归属，两个步骤不能混成一句。


\clearpage\section{最大可含锥：与基点无关}\label{g:109:114}

\textbf{对应原课件：109，第 114--118 页。}


\paragraph{先从衰退定义读出集合恒等式。}
非空闭凸 $C$ 有 $C+C_\infty=C$：左到右因为每个渐近方向是衰退方向、取步长1；右到左因为 $0\in C_\infty$，每点可写 $c=c+0$。若另有 $0\in C$，取 $c=0$ 立即得 $C_\infty\subseteq C$。下面还保留老师的截尾证明，说明同一结论怎样不用先引用衰退定理得到。
\paragraph{最大可含锥不是“随便选一个尽量大的锥”。}
固定非空闭凸 $S$ 中基点 $x_0$，所求是一个包含\emph{所有}合格锥的最大锥，合格条件为 $x_0+C'\subseteq S$。老师给出显式公式
\[C=\bigcap_{\lambda>0}\lambda(S-x_0).
\]
要证三个部分：此公式本身是合格锥且包含任何合格锥；换基点不改变它；它等于 $AS$。下面逐部分核验。$x_0$ 不是要求为原点，用 $S-x_0$ 是把它暂时平移到原点。


\subsection{命题与证明：原页 114}

\textbf{命题。}若闭凸集含0，则其渐近锥包含于它本身。

\paragraph{老师的证明思路。}把0代入已证 $C+C_\infty=C$。

\paragraph{逐步推导。}
假设 $C$ 闭凸且 $0\in C$。由上一命题 $C+C_\infty=C$。任取 $d\in C_\infty$，因 $0\in C$，和集定义给 $0+d\in C+C_\infty=C$。所以 $d\in C$，得到 $C_\infty\subseteq C$。
这是老师本页实际给的短推论，长度虽短但逻辑已全：明确选任意方向、使用0作为原集合点、代入已证和集恒等式，再返回包含结论。它无需另起一条新的完整衰退锥证明。


\subsection{命题与证明：原页 115}

\textbf{命题。}含0的闭凸集包含其渐近锥。

\paragraph{老师的证明思路。}将渐近方向见证点缩小到0与该点的线段上，最后用闭性。

\paragraph{逐步推导。}
老师这一版使用截尾生成闭锥，须补足闭包逼近。设闭凸 $C$ 含0，$d\in AC$。若 $d=0$，已在 $C$。否则选足够大的阈值整数 $M$。由 $d\in\Gamma(C^M)$，可取 $z_M\in C$、$\|z_M\|\ge M$ 和 $\alpha_M\ge0$ 使
\[\|\alpha_Mz_M-d\|<1/M.\]
当 $M>\|d\|+1$ 时，$\alpha_M\|z_M\|\le\|d\|+1/M$，所以
\[0\le\alpha_M\le(\|d\|+1/M)/M<1.\]
因此 $\alpha_Mz_M=(1-\alpha_M)0+\alpha_Mz_M$ 是两个 $C$ 中端点的凸组合，在 $C$ 中。又选点误差使这些组合趋于 $d$，闭性给 $d\in C$。
原页直接说存在一个缩放点恰等于 $d$，忽略了生成\emph{闭}锥可能只含缩放极限；加上逼近与最后闭性才能严谨。保留的核心思路仍是“沿0与远点之间线段缩回”，并没有换成别的分离定理。


\subsection{命题与证明：原页 117、118}

\textbf{命题。}闭凸集在某基点可容纳的最大锥等于 $\bigcap_{\lambda>0}\lambda(C-a)$，与基点无关，并等于渐近锥。

\paragraph{老师的证明思路。}先把交归属变为整条射线归属，再借凸性和闭性将射线平移到任意另一基点。

\paragraph{逐步推导。}
设非空闭凸 $S$，固定 $a\in S$，令 $D_a=\bigcap_{\lambda>0}\lambda(S-a)$。老师先描述 $a+D$ 能完全留在 $S$ 中的最大锥，再给此交表达式。
\paragraph{第一部分：交等于从 $a$ 出发的全部可行射线方向}
对任意 $d$，$d\in D_a$ 等价于对每个 $\lambda>0$，$d/\lambda\in S-a$，即 $a+d/\lambda\in S$。以 $t=1/\lambda$ 重命名，正好是 $a+td\in S$ 对所有 $t>0$；$t=0$ 由 $a\in S$ 补上。这个方向集含0，且非负倍数仍保持整条射线，所以确为锥。任何 $a+C'\subseteq S$ 的锥中方向 $d$ 都满足 $a+td\in S$，因而 $C'\subseteq D_a$；而 $a+D_a\subseteq S$，所以它就是最大锥。原页反向直接从点在交中说整射线在交中，缺的正是这个对任意 $\lambda$ 的变量替换。
\paragraph{第二部分：基点无关，补足原页未经论证的等价}
给 $d\in D_a$ 与任意另一基点 $b\in S$，要证对每个固定 $t>0$ 都有 $b+td\in S$。取 $T\to\infty$ 且 $T\ge t$，由原射线 $a+Td\in S$ 及凸性，
\[z_T=(1-t/T)b+(t/T)(a+Td)\in S.\]
展开 $z_T=b+td+(t/T)(a-b)\to b+td$，闭性给极限在 $S$。因此 $d\in D_b$；交换 $a,b$ 得反向，故 $D_a=D_b$。原页只写 $a+\R_+d\subset S\iff b+\R_+d\subset S$ 而没有理由，这段凸组合极限就是不可省略的证明。
\paragraph{第三部分：与渐近锥一致}
若 $d\in D_a$，取 $x_n=a+nd\in S,t_n=n$，则 $x_n/t_n\to d$，故 $D_a\subseteq S_\infty$。反向给 $d\in S_\infty$，取 $x_n/t_n\to d$，对任意固定 $t>0$ 以 $(1-t/t_n)a+(t/t_n)x_n$ 作凸组合，取极限得到 $a+td\in S$，故 $d\in D_a$。这个步骤正是截尾与射线之间的逻辑桥梁，不能把“在生成闭锥中”直接当成“原集合包含整条射线”。


\clearpage\section{109自测：把归一化和凸组合练熟}
\subsection{计算三个渐近锥}
求 $[2,5]$、$[2,\infty)$ 和 $C=\R\times[0,1]$ 的渐近锥。
\paragraph{解答。}有界非空集 $[2,5]$ 的锥是 $\{0\}$。对 $[2,\infty)$，比值非负使负方向不可能；任给 $d\ge0$ 取 $x_n=2+nd,t_n=n$，得到方向 $d$，故锥 $[0,\infty)$。对条带，任意见证 $(a_n,b_n)/t_n$ 中 $0\le b_n\le1$，所以第二坐标趋于0。任意 $(d,0)$ 可用 $(nd,0)\in C$、$t_n=n$ 实现，所以 $C_\infty=\R\times\{0\}$。
\subsection{闭集加紧集为什么闭}
设 $A$ 闭，$B$ 紧且非空，证明 $A+B$ 闭，并说明与正半独立定理的关系。
\paragraph{解答。}取 $z_n=a_n+b_n\to z$。$B$ 紧给 $b_{n_k}\to b\in B$，从而 $a_{n_k}=z_{n_k}-b_{n_k}\to z-b$。$A$ 闭使 $z-b\in A$，所以 $z\in A+B$。锥语言中 $B_\infty=\{0\}$，因此 $d_A+d_B=0$ 自动迫使二者为0；两个锥正半独立，闭和定理也适用。这里直接证明更短，但两种机制一致。
\subsection{正半独立条件怎样检查}
$A=[0,\infty)$、$B=(-\infty,0]$ 的渐近锥是否正半独立？它们的和是否闭？
\paragraph{解答。}它们本来就是闭锥，故渐近锥分别等于自身。取 $1\in A_\infty,-1\in B_\infty$ 相加为0却都非零，所以不正半独立。但 $A+B=\R$ 仍闭。这说明定理给的是充分条件；条件失败不能反推和必不闭。
\subsection{复写归一化矛盾}
在 $\ker L\cap C_\infty=\{0\}$ 下，证明 $x_n\in C$ 且 $(Lx_n)$ 有界必有 $(x_n)$ 有界。
\paragraph{解答。}若无界，可抽子列使 $r_n=\|x_n\|\to\infty$；单位列 $u_n=x_n/r_n$ 有子列趋于 $u$ 且 $\|u\|=1$。由同一 $x_n,r_n$ 见证 $u\in C_\infty$；线性连续性给 $Lu=\lim Lx_n/r_n=0$。核条件便使 $u=0$，与单位范数矛盾。三项不能漏：标度趋无穷、单位范数保非零、像有界使商趋零。
\subsection{闭凸性各用在何处}
设非空闭凸 $C$、$d\in C_\infty$。证明对任意 $a\in C,t\ge0$，$a+td\in C$。
\paragraph{解答。}选 $x_n\in C,s_n\to\infty,x_n/s_n\to d$。$t=0$ 直接成立；$t>0$ 时充分大 $s_n\ge t$，比例 $t/s_n\in[0,1]$。因此凸性给 $(1-t/s_n)a+(t/s_n)x_n\in C$。此列趋于 $a+td$，闭性才让极限仍在 $C$。基点 $a$ 和长度 $t$ 都任意，所以得衰退方向。凸性用来制造集内点，闭性用来留下极限。
\subsection{凸包和的四项展开}
设 $a=\tfrac13a_1+\tfrac23a_2$、$b=\tfrac14b_1+\tfrac34b_2$，把 $a+b$ 写成 $A+B$ 中点的凸组合。
\paragraph{解答。}
\[a+b=\tfrac1{12}(a_1+b_1)+\tfrac14(a_1+b_2)
+\tfrac16(a_2+b_1)+\tfrac12(a_2+b_2).
\]
四个权重非负且和1。$a_1$ 的合并权重 $1/12+1/4=1/3$，$a_2$ 的合并权重 $1/6+1/2=2/3$；$b_1,b_2$ 类似。逐个核对系数才能确认得到的点确实为 $a+b$。
\subsection{计算一个极锥}
令 $C=\{(t,0):t\ge0\}$，求 $C^\circ$。
\paragraph{解答。}极锥条件为 $p_1t+p_2\cdot0\le0$ 对所有 $t\ge0$ 成立。取 $t=1$ 得 $p_1\le0$；反向 $p_1\le0$ 足以使所有这些不等式成立，$p_2$ 不受限制。所以 $C^\circ=\{(p_1,p_2):p_1\le0\}$。这是一个半空间，而不是仅仅朝左的射线。
\subsection{为什么渐近正则性有用}
如果 $d_1,d_2\in C_\infty$，能否直接把它们的两条见证相加来证明 $C_\infty$ 凸？
\paragraph{解答。}不能直接做，因为两条见证可能是 $x_n/s_n\to d_1$ 和 $y_n/t_n\to d_2$，标度不同。非空凸集的渐近正则性允许改用共同标度 $n$，分别选 $u_n/n\to d_1,v_n/n\to d_2$。凸性给 $z_n=\lambda u_n+(1-\lambda)v_n\in C$，于是 $z_n/n\to\lambda d_1+(1-\lambda)d_2$，成为合法的单一见证。先统一尺度再组合，是这一组定理的逻辑顺序。


\clearpage\section{原课件逐页对应索引}

下表每行是一张教师原课件页。续页、重复页及图示页随所在专题合并解释；新讲义页码指该专题的开始位置，具体证明另在正文标出原页。

\begin{longtable}{r p{.66\linewidth} r}\toprule 原课件页&本讲义对应专题&本讲义页\\\midrule\endhead

1 & 向量运算、连续性与平移 & \pageref{g:109:1}\\

2 & 向量运算、连续性与平移 & \pageref{g:109:1}\\

3 & 向量运算、连续性与平移 & \pageref{g:109:1}\\

4 & 向量运算、连续性与平移 & \pageref{g:109:1}\\

5 & 向量运算、连续性与平移 & \pageref{g:109:1}\\

6 & 向量运算、连续性与平移 & \pageref{g:109:1}\\

7 & 向量运算、连续性与平移 & \pageref{g:109:1}\\

8 & 集合和、集合差与闭包 & \pageref{g:109:8}\\

9 & 集合和、集合差与闭包 & \pageref{g:109:8}\\

10 & 集合和、集合差与闭包 & \pageref{g:109:8}\\

11 & 紧集的和：老师的两条证明 & \pageref{g:109:11}\\

12 & 紧集的和：老师的两条证明 & \pageref{g:109:11}\\

13 & 紧集的和：老师的两条证明 & \pageref{g:109:11}\\

14 & 紧集的和：老师的两条证明 & \pageref{g:109:11}\\

15 & 函数相加与对应相加 & \pageref{g:109:15}\\

16 & 函数相加与对应相加 & \pageref{g:109:15}\\

17 & 函数相加与对应相加 & \pageref{g:109:15}\\

18 & 线段、射线、锥与正半独立 & \pageref{g:109:18}\\

19 & 线段、射线、锥与正半独立 & \pageref{g:109:18}\\

20 & 线段、射线、锥与正半独立 & \pageref{g:109:18}\\

21 & 线段、射线、锥与正半独立 & \pageref{g:109:18}\\

22 & 线段、射线、锥与正半独立 & \pageref{g:109:18}\\

23 & 线段、射线、锥与正半独立 & \pageref{g:109:18}\\

24 & 三种无穷远描述：先把定义分开 & \pageref{g:109:24}\\

25 & 三种无穷远描述：先把定义分开 & \pageref{g:109:24}\\

26 & 三种无穷远描述：先把定义分开 & \pageref{g:109:24}\\

27 & 三种无穷远描述：先把定义分开 & \pageref{g:109:24}\\

28 & 三种无穷远描述：先把定义分开 & \pageref{g:109:24}\\

29 & 有界判据、平移不变与乘积包含 & \pageref{g:109:29}\\

30 & 有界判据、平移不变与乘积包含 & \pageref{g:109:29}\\

31 & 有界判据、平移不变与乘积包含 & \pageref{g:109:29}\\

32 & 有界判据、平移不变与乘积包含 & \pageref{g:109:29}\\

33 & 有界判据、平移不变与乘积包含 & \pageref{g:109:29}\\

34 & 渐近锥怎样控制交集的有界性 & \pageref{g:109:34}\\

35 & 渐近锥怎样控制交集的有界性 & \pageref{g:109:34}\\

36 & 渐近锥怎样控制交集的有界性 & \pageref{g:109:34}\\

37 & 线性像定理：归一化反证的完整链条 & \pageref{g:109:37}\\

38 & 线性像定理：归一化反证的完整链条 & \pageref{g:109:37}\\

39 & 线性像定理：归一化反证的完整链条 & \pageref{g:109:37}\\

40 & 线性像定理：归一化反证的完整链条 & \pageref{g:109:37}\\

41 & 线性像定理：归一化反证的完整链条 & \pageref{g:109:37}\\

42 & 闭集相加：正半独立排除抵消 & \pageref{g:109:42}\\

43 & 闭集相加：正半独立排除抵消 & \pageref{g:109:42}\\

44 & 闭集相加：正半独立排除抵消 & \pageref{g:109:42}\\

45 & 闭集相加：正半独立排除抵消 & \pageref{g:109:42}\\

46 & 凸集的定义、运算与连通性 & \pageref{g:109:46}\\

47 & 凸集的定义、运算与连通性 & \pageref{g:109:46}\\

48 & 凸集的定义、运算与连通性 & \pageref{g:109:46}\\

49 & 凸集的定义、运算与连通性 & \pageref{g:109:46}\\

50 & 凸集的定义、运算与连通性 & \pageref{g:109:46}\\

51 & 凸集的定义、运算与连通性 & \pageref{g:109:46}\\

52 & 凸集的定义、运算与连通性 & \pageref{g:109:46}\\

53 & 凸集的定义、运算与连通性 & \pageref{g:109:46}\\

54 & 凸集的定义、运算与连通性 & \pageref{g:109:46}\\

55 & 凸包、闭凸包与凸包和公式 & \pageref{g:109:55}\\

56 & 凸包、闭凸包与凸包和公式 & \pageref{g:109:55}\\

57 & 凸包、闭凸包与凸包和公式 & \pageref{g:109:55}\\

58 & 凸包、闭凸包与凸包和公式 & \pageref{g:109:55}\\

59 & 凸包、闭凸包与凸包和公式 & \pageref{g:109:55}\\

60 & 内积、超平面与 Minkowski 定理 & \pageref{g:109:60}\\

61 & 内积、超平面与 Minkowski 定理 & \pageref{g:109:60}\\

62 & 内积、超平面与 Minkowski 定理 & \pageref{g:109:60}\\

63 & 内积、超平面与 Minkowski 定理 & \pageref{g:109:60}\\

64 & 内积、超平面与 Minkowski 定理 & \pageref{g:109:60}\\

65 & 极锥与坐标偏序 & \pageref{g:109:65}\\

66 & 极锥与坐标偏序 & \pageref{g:109:65}\\

67 & 重新审视渐近方向与衰退方向 & \pageref{g:109:67}\\

68 & 重新审视渐近方向与衰退方向 & \pageref{g:109:67}\\

69 & 重新审视渐近方向与衰退方向 & \pageref{g:109:67}\\

70 & 重新审视渐近方向与衰退方向 & \pageref{g:109:67}\\

71 & 重新审视渐近方向与衰退方向 & \pageref{g:109:67}\\

72 & 重新审视渐近方向与衰退方向 & \pageref{g:109:67}\\

73 & 重新审视渐近方向与衰退方向 & \pageref{g:109:67}\\

74 & 重新审视渐近方向与衰退方向 & \pageref{g:109:67}\\

75 & 重新审视渐近方向与衰退方向 & \pageref{g:109:67}\\

76 & 渐近锥闭性、闭包不变与锥的特例 & \pageref{g:109:76}\\

77 & 渐近锥闭性、闭包不变与锥的特例 & \pageref{g:109:76}\\

78 & 渐近锥闭性、闭包不变与锥的特例 & \pageref{g:109:76}\\

79 & 渐近锥闭性、闭包不变与锥的特例 & \pageref{g:109:76}\\

80 & 渐近锥闭性、闭包不变与锥的特例 & \pageref{g:109:76}\\

81 & 三个渐近描述为什么等价 & \pageref{g:109:81}\\

82 & 三个渐近描述为什么等价 & \pageref{g:109:81}\\

83 & 三个渐近描述为什么等价 & \pageref{g:109:81}\\

84 & 用序列重证有界性与平移性质 & \pageref{g:109:84}\\

85 & 用序列重证有界性与平移性质 & \pageref{g:109:84}\\

86 & 用序列重证有界性与平移性质 & \pageref{g:109:84}\\

87 & 用序列重证有界性与平移性质 & \pageref{g:109:84}\\

88 & 乘积与交：不能交换的运算 & \pageref{g:109:88}\\

89 & 乘积与交：不能交换的运算 & \pageref{g:109:88}\\

90 & 乘积与交：不能交换的运算 & \pageref{g:109:88}\\

91 & 乘积与交：不能交换的运算 & \pageref{g:109:88}\\

92 & 再走一遍线性像定理 & \pageref{g:109:92}\\

93 & 再走一遍线性像定理 & \pageref{g:109:92}\\

94 & 再走一遍线性像定理 & \pageref{g:109:92}\\

95 & 再走一遍线性像定理 & \pageref{g:109:92}\\

96 & 再走一遍线性像定理 & \pageref{g:109:92}\\

97 & 闭和定理：线性映射法与直接子列法 & \pageref{g:109:97}\\

98 & 闭和定理：线性映射法与直接子列法 & \pageref{g:109:97}\\

99 & 闭和定理：线性映射法与直接子列法 & \pageref{g:109:97}\\

100 & 闭和定理：线性映射法与直接子列法 & \pageref{g:109:97}\\

101 & 闭和定理：线性映射法与直接子列法 & \pageref{g:109:97}\\

102 & 闭和定理：线性映射法与直接子列法 & \pageref{g:109:97}\\

103 & 闭和定理：线性映射法与直接子列法 & \pageref{g:109:97}\\

104 & 渐近正则：把不同尺度统一 & \pageref{g:109:104}\\

105 & 渐近正则：把不同尺度统一 & \pageref{g:109:104}\\

106 & 渐近正则：把不同尺度统一 & \pageref{g:109:104}\\

107 & 渐近正则：把不同尺度统一 & \pageref{g:109:104}\\

108 & 渐近正则：把不同尺度统一 & \pageref{g:109:104}\\

109 & 闭凸集的渐近锥等于衰退锥 & \pageref{g:109:109}\\

110 & 闭凸集的渐近锥等于衰退锥 & \pageref{g:109:109}\\

111 & 闭凸集的渐近锥等于衰退锥 & \pageref{g:109:109}\\

112 & 闭凸集的渐近锥等于衰退锥 & \pageref{g:109:109}\\

113 & 闭凸集的渐近锥等于衰退锥 & \pageref{g:109:109}\\

114 & 最大可含锥：与基点无关 & \pageref{g:109:114}\\

115 & 最大可含锥：与基点无关 & \pageref{g:109:114}\\

116 & 最大可含锥：与基点无关 & \pageref{g:109:114}\\

117 & 最大可含锥：与基点无关 & \pageref{g:109:114}\\

118 & 最大可含锥：与基点无关 & \pageref{g:109:114}\\

\bottomrule\end{longtable}

\paragraph{材料与记号说明。}按用户提供的教师原课件 109-ZSM-PPT.pdf 重新编写；相关基础背景参照 Debreu 的 Theory of Value。正文以教师原页为定位，不以旧 gpt 笔记分页为定位。补充证明与原稿纠正均显式说明。